% Gerado por tools/tei2latex.py a partir de entrada.xml. Regenerar, não editar.

\pg{2} : * \textbar{} Elementos de philosophia, Thiago \textbar{} \emph{M}. Pae wx BM002777

\pg{3} BISPO TITULAR DE TIBERIADES

DOUTOR EM PHILOSOPHIA, THEOLOGIA E DIREITO CANONICO

MEMBRO DA ACADEMIA ROMANA DE 8. THOMAZ D'AQUINO SECRETARIO DA \&. C. DOS SEMINARIOS E DAS UNIVERSIDADES DOS ESTUDOS

\textbar{} \textbar{} \textbf{ELEMENTOS}

\textbf{PHILOSOPHIA}

\textbar{} Ec \textbf{cm}

\textbf{VOLUME II}

\textbf{ANTHROPOLOGIA} - \textbf{THEODICEA} - \textbf{MORAL}

QUARTA EDICÁO 4

ROMA

"Via DEL Banco 8. SPrIRITO, 12

1916

\pg{4} \textbf{APPROBATIO}

Nova mandetur typis editio philosophiae lusitano idiomate exaratae, cuius inscriptio est: \emph{Elementos de philosophia}, auctore Reverendissimo et Amplissimo Domino lacobo Sinibaldi, episcopo Tiberiadis.

\emph{Assisii}, \emph{die} 30 \emph{maii} 19106.

T FR. AMBROSIUS EPs.

Assisi --- Tipografia Metastasio --- Assisi

\pg{5} \textbf{METAPHYSICA}

%%>>> parte-1

\part*{PARTE TERCEIRA\\[0.3em]\large ANTHROPOLOGIA}
\addcontentsline{toc}{part}{PARTE TERCEIRA --- ANTHROPOLOGIA}

PROLOGO

%%>>> t1

\tese{1}{Anthropologia.}
\emph{Anthropologia} É a \emph{sciencia que trata das causas supremas do homem}. --- O \emph{homem} define -\emph{se:.animal racional} \footnote{\emph{Animal racional} é a definição classica do \emph{homem}, a qual, desde a mais remota antiguidade, foi recebida pelas maiores intelligencias. --- Esta definição é perfeita, porque \emph{animal} é o \emph{genero proximo} do homem, e \emph{racional} é a sua \emph{differença especifica}. Emquanto \emph{animal}, o homem é similhante aos seres sensitivos; emquanto \emph{racional}, é similhante aos seres meramente espirituaes. D'este modo o homem é o annel de conjuncção entre o ceu e a terra, entre os Anjos e os animaes irracionaes. s\par Pergunta-se: Porque é que o \emph{tomem} se define: \emph{animal racional}, e não: \emph{animal intelligente} ?\par Resposta. A \emph{intelligencia} e a \emph{razão} não são duas faculdades distinctas, mas indicam duas diversas funcções da mesma faculdade. Na verdade, a \emph{intelligencia}, emquanto por uma \emph{simples intuição} percebe o \emph{immaterial}, o \textbar{} \emph{universal}, o \emph{abstracto}, conserva o nome de \emph{tintelligencia}; --- mas, emquanto \emph{raciocina}, isto é, emquanto de um principio universal deduz consequencias particulares, toma, em \emph{sentido mais rigoroso}, o, nome de \emph{razão}. Ora, sendo proprio do homem conhecer a verdade, não \emph{immediatamente} e pela simples \emph{intuição}, mas \emph{mediatamente} e pelo \emph{raciocinio}, é claro que o homem deve \textbar{} dizer-se animal \emph{racional}, e não animal \emph{intelligente}. \textbar{}\par Dizemos --- em \emph{sentido mais rigoroso}; porque, em sentido mais \emph{lato}, a razão significa tambem a faculdade que percebe, por \emph{simples éiutwicdo}, o \textbar{} \emph{immalterial}, o \emph{universal}, o \emph{abstracto} (e então diz-se razão \emph{intuitiva}); e neste \textbar{} sentido costuma attribuir-se tambem o \textsc{DEus} e aos Anjos, que conhecem a verdade não pelo \emph{raciocinio}, mas pela \emph{intuição}. ; ;\par \emph{Antonio Rosmini} diz que esta definição de \emph{homem} é defeituosa por tres razões: --- 1.º) porque a differenca \emph{racional} denota a \emph{intelligencia} do homem, mas não a \emph{vontade}; --- 92.) porque a \emph{razão} deriva da \emph{intelligencia} e numa definição deve apontar-se a causa, e não o effeito; - 3.º) porque esta \textbar{}\pgn{6}definição --- \emph{animal racional} --- parece indicar que no homem o sujeito é o \emph{animal} e que a \emph{razão} é apenas um \emph{attributo} ou uma \emph{propriedade}, que se Ihe \&ecrescenta, depois de constituido na sua essencia.\par As tres razões, por que \emph{Rosmini} rejeita aquella definição do //\emph{omem}, são infundadas, --- E' infundada a \emph{primeira}; porque a \emph{vontade} é uma consequencia natural da \emph{razdo}, como o proprio \emph{Rosmini} reconhece, e nas definições basta apontar o elemento \emph{principal}, de que os outros derivam. --- E' infundada a \emph{segunda}; porque a \emph{intelligencia} e a \emph{vasdo} não são duas faculdades distinctas, mas são uma e a mesma faculdade, considerada em duas diversas funcções. --- E' infundada a \emph{terceira}; porque na definição do \emph{homem} a palavra « \emph{racional} » não denota um simples attributo do \emph{animal}, mas é a dif ferenca \emph{especifica}, que, unida ao genero, constitue a essencia do homem,}. '

%%>>> t2

\tese{2}{Objecto da Anthropologia.}
O \emph{objecto} da \emph{Anthropologia} é o \emph{homem}, considerado nas suas \emph{causas swupremas}. --- As causas \pg{6}  \emph{supremas} do homem são --- a sua \emph{essencia}, que se compõe da causa \emph{material} e da \emph{formal}, e de que derivam as \emph{faculdades} e. as \emph{operações}, --- a sua \emph{origem}, isto é, á sua causa \emph{efficiente}, --- e o seu \emph{destino}, isto é, a sua causa \emph{final}.

D'aqui se vê a differença entre a \emph{Anthropologia} e a \emph{Physiologia}, que tambem tem por objecto o \emph{homem}; pois a \emph{Physiologia} limita-se a considerar e classificar as funcções do organismo humano, ao passo que a \emph{Anthropologia}, baseada nos phenomenos que se realizam no homem, chega ao conhecimento dos mais altos principios, das causas supremas, de que os phenomenos derivam \footnote{Os \emph{phenomenos} ou \emph{factos}, que se realizam no homem, podem reduzir-se a tres grupos principaes. Uns referem-se á vida \emph{vegetativa}, --- outros, á yida \emph{sensitiva}, --- outros, \& vida \emph{intellectual}. :\par Os phenomenos, que se referem á vida \emph{vegetativa} e \emph{sensitiva}, chamam-se \emph{physiologicos}, e constituem o objecto da Physiologia; --- os phenomenos \emph{intellectuaes} dizem-se \emph{psychologicos}, e d'elles se occupa a Psychologia, que é uma parte da Anthropologia.\par Os phenomenos \emph{psychologicos} e os \emph{physiologicos} differem entre si principalmente --- a) na sua \emph{natureza}, --- \emph{b}) \emph{escopo}, --- \emph{c}) e \emph{cognoscibilidade}.\par a) Differem na sua natureza. Os phenomenos \emph{psychologicos} são \emph{espirituaes} (como os \emph{pensamentos} o as \emph{volíções}), não têm por si extensão nem figura, não se podem dividir nem medir \emph{directamente}. Os phenomenos \emph{physiologicos} são \emph{materiaes}, occupam um logar no espaço e no tempo e são divisiveis e commensuraveis.\par \emph{b}) Differem no \emph{escopo}. Os phenomenos \emph{psychologicos} tem por fim a percepção e o amor de um objecto \emph{espiritual}. Os phenomenos \emph{Pphlysiologicos} tendem para a conservação do organismo ou para um objecto \emph{sensivel}.\par c) Differem na \emph{cognoscibilidade}. Os phenomenos \emph{psychologicos} conhecem-se por uma faculdade \emph{espiritual}. Os phenomenos \emph{physiologicos} ou são \emph{inconscientes} (como os da vida \emph{vegetativa}), ou conhecem-se por uma faculdade \emph{sensitiva} (como os phenomenos da vida \emph{sensitiva}).\par Todavia os phenomenos \emph{psychologicos} e os \emph{plysiologícos}, ainda que diffiram sob muitos respeitos, conservam entre si as seguintes relações:\par a) Todos derivam de um unico principio, que é \emph{a alma humana};\pgn{7}porque, como provaremos mais tarde, a \emph{alma humana} é forma \emph{substancial} do corpo, e a forma é o unico principio de toda a actividade do composto.\par \emph{b}) Os phenomenos \emph{psychologicos} são sempre acompanhados por phenomenos \emph{physiologicos}; e isto verifica-se nos pensamentos e nas volições, que, embora deriveni unica e exclusivamante da alma espiritual, todavia estão ligados ás funcções \emph{organicas}, não como a \emph{causas}, mas como a \emph{condições} indispensáveis para a sua producção.\par D'ahi se vê que a \emph{Anthropologia} e a \emph{Physiologia}, posto que sejam duas sciencias distinctas, porque os seus objectos são distinctos, ainda assim se relacionam e se prestam mutuo auxilio.}.

%%>>> t3

\pg{7} \tese{3}{Excellencia da Anthropologia.}
A excellencia da \emph{Anthropologia} demonstra-se pela excellencia do seu \emph{objecto}. Por quanto, o \emph{Homem} não só concentra em si as perfeicoes de todas as outras creaturas corporeas, que lhe estão sujeitas, mas tambem participa dos dotes das intelligencias separadas, que lhe estão superiores; pois que possue o \emph{ser} como os \emph{mineraes}, \emph{vive} como os \emph{vegetaes}, \emph{sente} como os \emph{anmaes}, \emph{entende} como os \emph{espiritos}. D' este modo o homem é annel de conjuneção entre a materia e o espirito, e com razão chama-se \emph{microcosmos} (pequeno mundo) \footnote{O \emph{Homem} é um \emph{pequeno mundo} --- não só pela sua natureza, emquanto é o resumo, a synthese de todas as perfeições creadas, --- mas tambem pelo seu \emph{conhecimento}, emquanto reproduz e exprime na sua intelligencia o mundo corporeo e o espiritual, e concebe, embora de um modo imperfeito e inadequado, a idea de um Creador sapientissimo e omnipotente,}.

%%>>> t4

\tese{4}{Methodo da Anthropologia.}
O methodo da \emph{Anthropologia} --- \emph{a}) não deve ser exclusivamente \emph{analytico}, --- 0) nem exclusivamente \emph{synthetico}, --- \emph{c}) mas deve ser \emph{analytico-synthetico}.

a) Não deve ser exclusivamente \emph{analytico}; pois que a \emph{experiencia}, de que este methodo parte, não pode, só por si, attingir a essencia da alma e pode levar ao materialismo.

b) Nem deve ser exclusivamente \emph{synthetico}, isto é, fundado unicamente nos \emph{principios} da \emph{razão}; porque a razão, sem o auxilio da experiencia, forma hypotheses arbitrarias, que não correspondem á realidade, e pode levar ao pantheismo e a outros erros perniciosos.

c) Mas deve ser \emph{analytico-synthetico}. --- Pelo methodo \emph{analytico}, \emph{observamos} primeiramente os phenomenos, tanto \emph{internos} como \emph{externos}, considerando-os em si e em todas as suas condições e eireumstaneias; depois, em virtude do \emph{principio de causalidade}, \emph{subimos} ás causas \emph{proximas} e \emph{immediatas} d'aquelles phenomenos, e d'estas para a \emph{mais elevada} e \emph{suprema} \pg{8} , que é a \emph{essencia} do homem. --- Pelo \emph{synthetico}, fazemos o caminho inverso; partimos da \emph{essencia} do homem, já conhecida, e descemos ás suas \emph{faculdades}, e d'estas ás \emph{operações}.

Em menos palavras: Na \emph{Anthropologia} deve empregar-se a \emph{inducção} e o \emph{syllogismo}. --- Pela \emph{inducção}, \emph{subimos} do effeito á causa, e assim conhecemos a causa pelo effeito. --- Pelo \emph{syllogismo}, \emph{descemos} da causa ao affeito, e assim conhecemos o effeito pela causa \footnote{A primeira \emph{funcção} do methodo \emph{analytico} é a \emph{observação}, tanto \emph{interna} come \emph{externa}. --- A. observação \emph{interna} costuma chamar-se \emph{methodo subjectivo}; a a\&bservacüo \emph{externa} diz-se \emph{methodo objectivo}. --- Digamos alguma coisa ácerca d'estes dois methodos.\par a) O methodo \emph{subjectivo} consiste em observar as operações ou modificações \emph{internas}, por meio da \emph{consciencia} ou percepção \emph{interna}. --- Na \emph{Logica} demonstrámos que a \emph{consciencia} não se engana na percepção do seu objecto.\par O methodo \emph{subjectivo} é necessario na \emph{Anthropologia}. Na verdade, a existencia e a natureza da alma não se conhecem directamente em si, mas 8ó nas operações da vida sensitiva e intellectual, de que á mesma alma é principio. Ora as operações da vida sensitiva e intellectual realizam-se no \emph{intimo} do homem, e por isso só podem ser conhecidas pelo methodo \emph{subjectivo}.\par Este methodo serve-se tambem da \emph{experimentação}.\par Alguns escriptores só admittem a utilidade da \emph{experimentação} nas sciencias \emph{naturaes}, e dizem que a \emph{experimentação} na \emph{Anthropologia} é impossivel.\par Esta opinião não póde acceitar-se. Confessamos de boa mente que a \emph{experimentação} na \emph{Anthropologia} não é \emph{sempre} possivel e que muitas vezes é incerta nos resultados; mas sustentamos que não é ella de todo e sempre impossivel, porque podemos reproduzir um phenomeno interno, como um \emph{pensamento}, um \emph{raciocinio}, uma \emph{volição}, etc., para estudar a natureza d'esse acto e o modo por que se produz.\par b) O methodo \emph{objectivo} consiste em conhecer e explicar os phenomenos \emph{psychologicos} por meio do estudo das diversas \emph{raças humanas}, da \emph{historia}, e das \emph{linguas}. --- Este methodo, empregado com discrição, completa e esclarece os dados do methodo \emph{subjectivo}.\par O methodo \emph{objectivo} funda-se principalmente no estudo da \emph{psychologia comparada}, da \emph{historia} e das \emph{linguas}.\par a) A \emph{psychologia comparada} observa a criança, em quem a razão ainda não se desenvolveu, --- o demente, em quem ella se eclypsou, --- o selvagem, em quem não se aperfeiçoou; compara o homem com o animal, e conhece quaes são as faculdades proprias de cada um, e quas são as communs a ambos. Estas observações projectam muita luz sobre a natureza e as leis das faculdades humanas.\par \emph{b}) \emph{A historia}, contemplando a humanidade inteira como um sô homem, reconhece nella as mesmas faculdades, os mesmos actos e sentimentos, que se encontram no individuo; e por isso o seu estudo é muito util para o conhecimento do homem, na ordem intellectual e moral.\par c) As \emph{linguas} reproduzem os pensamentos e as reflexões dos que fallam;\pgn{9}por isso o estudo da lingua de diversos povos ou de um só povo manifesta as diversas formas da vida intellectual e o seu desenvolvimento progressivo.\par Os escriptores não estão do accordo com relação ao methodo da \emph{Anthiropologia}, mas emittem e seguem diversas opiniões. Limitamo-nos a expôr e criticar as dos \emph{materialistas}, dos \emph{positivistas} e dos \emph{idealistas}.\par a) Os \emph{materialistas}, por identificarem os phenomenos \emph{psychologicos} com os \emph{physiologicos}, dizem que o methodo, que se ha de seguir na \emph{Ant/ropologia}, deve ser exclusivamente \emph{objectivo}, como na \emph{P/ysiología}.\par Esta opinião é falsa pelas seguintes razões: --- 1.º) porque identifica os phenomenos \emph{psycholoyicos} com os \emph{pysiologicos}, que, como vimos, são distinctos entre si; --- 2.º) porque o methodo \emph{objectivo} não póde prescindir do \emph{subjectivo}; pois que só pelo testemunho da percepção \emph{interna} podemos ficar certos de que um phenomeno foi observado pela percepção \emph{externa} com todas as precauções e com o emprego de todos os meios legitimos.\par b) Os \emph{positivistas} não negam nem affirmam a existencia da alma; dizem que a alma é \emph{zncognoscivel}, e que por isso a \emph{Anthropologia} deve limitar-se a descrever o systema nervoso e o mechanismo das acções moleculares, e empregar o methodo exclusivamente \emph{objectivo}.\par Esta opinião tambem é falsa pelas razões que acima adduzimos contra os materialistas. --- Só aecrescentaremos que, quando os \emph{positivistas} affirmam que a nossa alma é 2\emph{ncognoscivel}, a sua affirmação é \emph{falsa} ou \emph{gratuita}. Na verdade, ou elles conhecem a alma humana, ou não a conhecem. Se a conhecem, é \emph{falso} que ella seja \emph{Zncognoscivel}. Se não a conhecem, \emph{gratuitamente} affirmam que ella é \emph{éncognoscivel}.\par c) Os \emph{idealistas}, partindo do principio --- de que todas as coisas são ideas, ás quaes nenhuma realidade corresponde, --- affirmam que o methodo da \emph{Anthropologia} deve ser exclusivamente \emph{synthetico}, isto é, unicamente baseado nos principios da razão.\par Tambem esta opinião não póde acceitar-se, porque se funda num falso supposto. Às coisas são o que são, independentemente do nosso pensamento; e, se quizermos conhecel-as como são em si mesmas, é necessario servirmo-nos da experencia, tanto interna como externa.\par Coneluimos que o verdadeiro methodo, que se deve seguir e que seguiremos na \emph{Antropologia}, é o \emph{analytico-synthetico},}.

%%>>> t5

\pg{9} \tese{5}{Divisão da Anthropologia.}
Dividimos a \emph{Anthropologia} em \emph{quatro} secções. Na \emph{primiera} secção ocenpamo-nos da \emph{essencia} do homem; --- na \emph{segunda}, das suas \emph{faculdades}; --- na \emph{terceira}, das suas \emph{operações}; --- na \emph{quarta}, da sua \emph{origem} e do seu \emph{destino}.

10 METAPHYSICA --- PARTE TERCEIRA --- ANTHROPOLOGIA

\part*{SECÇÃO PRIMEIRA\\[0.3em]\large AA}
\addcontentsline{toc}{part}{SECÇÃO PRIMEIRA --- AA}

\pg{10} \textbf{ESSENCIA DO HOMEM}

%%>>> t6

\tese{6}{Divisão da secção primeira.}
Como a \emph{essencia} do homem é composta de alma e de corpo, substancial e pessoalmente unidos entre si, dividimos a secção \emph{primeira} em \emph{tres} capitulos. No \emph{primeiro} capitulo tratamos do \emph{corpo}; no \emph{segundo}, da \emph{alma}; no \emph{terceiro}, da \emph{união} da alma com o corpo.

%%>>> cap-1

\chapter*{CAPITULO PRIMEIRO\\[0.3em]\large CORPO HUMANO}
\addcontentsline{toc}{chapter}{CAPITULO PRIMEIRO --- CORPO HUMANO}

\textsc{SUNMMARIO}: --- Descripção do corpo humano. --- Orgãos da vida vegetativa.

--- Orgãos da sensibilidade externa. --- Orgãos da sensibilidade interna

e da locomoção.

%%>>> art-1

\section*{Artigo I.\\ Descripção do corpo humano.}
\addcontentsline{toc}{section}{Artigo I. --- Descripção do corpo humano.}

%%>>> t7

\tese{7}{Principaes regiões do corpo humano.}
As principaes regioes do corpo humano são \emph{tres}, a saber: a \emph{cabeça}, o \emph{thoraz}, e o \emph{abdomen}. --- O \emph{thorax} e o \emph{abdomen}, que estão separados pelo \emph{diaphragma}, constituem o \emph{tronco}, á que se unem, em forma de appendices, os \emph{membros thoracicos} e \emph{abdominaes}. --- De cada uma d'estas regiões faremos uma breve deseripção \footnote{Os elementos figurados do corpo humano, que é o mais perfeito e maravilhoso de todos os organismos, são compostos, em ultima analyse, de \emph{cellulas}, constituidas essencialmente por um protoplasma, capaz de se nuirir e reproduzir. --- As \emph{cellulas} primitivas multiplicam-se, tomam uma disposição diversa segundo a diversidade da sua natureza, e concorrem para formar os \emph{tecidos}, que geralmente se reduzem a seis typos principaes, a saber, o \emph{cellular} ou \emph{fibroso}, o \emph{conjuncirvo}, o \emph{muscular}, o nervoso, o \emph{osseo}, e o \emph{cartilagineo}. --- A reunião dos tecidos forma os \emph{orgãos}, que são instrumentos capazes de produzir certas e determinadas \emph{funcções}. Quando para ó exercicio de uma \emph{funcção} são necessarios muitos e diversos orgãos, a reunião d'esses orgãos diz-se \emph{apparelho}.}.

%%>>> t8

\tese{8}{Cabeça.}
A \emph{cabeça}, que é a região mais nobre do corpo humano, divide-se em. duas partes: o \emph{crameo} e a \emph{face}.

a) O \emph{craneo} é uma \emph{caixa ossea}, de forma oval, que aloja e protege o \emph{cerebro}, o \emph{cerebello} e a \emph{medulla allongada}. --- As \emph{paredes} superiores, lateraes e posteriores do \emph{craneo} são cobertas de pelle; a sua parede inferior, que é a \emph{base}, articula-se sobre a primeira das \emph{verftebras}, que por isso se chama \emph{atlas}.

b) A \emph{face} contém e protege os orgãos da \emph{vista}, do \emph{olfacto}, do \emph{gosto}. --- Os orgãos do \emph{ouvido} relacionam-se com o \emph{craneo} \footnote{A face \emph{superior} da \emph{cabeça} é lisa e regular, e forma a \emph{região epicranica}. A face \emph{posterior} é espheroidal e forma a \emph{região occipital}. As faces \emph{lateraes} dividem-se em tres regiões, que são: uma superior, que se chama \emph{região temporal}; outra inferior e posterior, que é a \emph{região auricular}; e a terceira inferior e anterior, que é a \emph{região parotidia}. A face \emph{anterior} diz-se simplesmente \emph{fuce}, e revela a perfeição, a nobreza o a origem do homem,}.

%%>>> t9

\pg{11} \tese{9}{Thorax.}
\emph{Thorax} é uma cavidade, limitada adiante pelo \emph{esterno}, dos lados pelas \emph{costellas} e pelos \emph{musculos intercostaes}, por traz pela \emph{columna vertebral}, por baixo por uma membrana muscular, chamada \emph{diaphragma}, de cima pelo \emph{pescoço}.

O \emph{thorax} subdivide-se em tres cavidades secundarias: uma \emph{media} e duas /\emph{ateraes}. --- A cavidade \emph{media} aloja o \emph{coração} e os principaes vasos que d' elle derivam, uma parte da \emph{trachéa} e dos \emph{bronchios}, o \emph{esophago}, \emph{á aorta thoracica}, O \emph{canal thoracico}, etc. --- As cavidades /\emph{ateraes} contêm os dois \emph{pulmões}.

%%>>> t10

\tese{10}{Abdomen.}
\emph{Abdomen} é uma cavidade mais espaçosa que o \emph{thorax}. --- W' limitado em cima pelo \emph{diaphragma}, em baixo pela \emph{bacia}, atraz pela \emph{columna vertebral}, adiante e dos lados pela \emph{pelle} e pelos \emph{musculos abdominaes}.

O \emph{abdomen} contém os principaes orgãos do apparelho \emph{digestivo}, que são o \emph{estomago}, o \emph{intestino}, o \emph{figado}, O \emph{pancreas}, ete.

%%>>> t11

\tese{11}{Membros.}
Os \emph{membros} são \emph{quatro}: dois \emph{thoracicos}, e dois \emph{abdominaes}.

\emph{a}) Cada membro \emph{thoracico} compõe-se de quatro partes; \emph{hombro}, \emph{braço}, \emph{ante-brago} e \emph{mão}.

5) Cada membro \emph{abdominal} compõe-se tambem de quatro partes: \emph{anca} ou \emph{quadril}, \emph{coxa}, \emph{perna} e pé.

%%>>> art-2

\section*{Artigo II.\\ Orgãos da vida vegetativa.}
\addcontentsline{toc}{section}{Artigo II. --- Orgãos da vida vegetativa.}

%%>>> t12

\tese{12}{Orgãos da vida vegetativa.}
Os orgãos da vida \emph{vegetativa}, de que tencionamos tratar, são os que servem para a \emph{nutrição} do homem. Estes orgãos são muitos e diversos, e \pg{12}  constituem tantos \emph{apparelhos}, quantos são os principaes actos da funeção \emph{nutritiva}. Ora esses actos são tres, a saber: a \emph{digestão} (com a \emph{absorpção intestinal}), a \emph{circulação} e a \emph{respiração}; pois que --- pela \emph{digestão}, os alimentos são elaborados, --- pela \emph{circulação}, são elles distribuidos, --- pela \emph{respiração}, o sangue, que os transporta, é depurado. Logo os orgãos da vida \emph{vegetativa} constituem tres \emph{apparelhos}: o \emph{digestivo}, o \emph{circulatorio} ,e o \emph{respiratorio} \footnote{Os muitos e diversos orgãos, de que é composto o corpo humano, tendem ou para conservar a vida do homem, ou para o relacionar com o mundo exterior --- Os que tendem para conservar a vida do homem são especialmente os orgãos da vida \emph{vegetativa}: e os que o relacionam com o mundo exterior são os orgãos da \emph{sensibilidade}, tanto externa como interna, e os da \emph{locomoção}. --- D' estas duas especies de orgãos resultam duas diversas especies de funcções-as da vida \emph{vegetativa} e as da vida de \emph{relação},}.

%%>>> t13

\tese{13}{Orgãos do apparelho digestivo.}
Os orgãos, que constituem o apparelho \emph{digestivo}, são a \emph{bocca}, a. \emph{pharinge}, o \emph{esophago}, o \emph{estomago}, o \emph{intestino delgado}, o \emph{intestino grosso}.

a) \textbf{Bocca}, --- A \emph{bocca}, primeira parte do canal \emph{digestivo}, é uma cavidade oval, situada na face entre as maxillas, limitada adiante pelos \emph{labios}, atraz pelo \emph{veo palatino}, dos lados pelas \emph{faces}, em baixo pela \emph{lingua}, e em cimá pela \emph{abobada palatina}. --- As \emph{mazcillas}, situadas uma por cimá da outra, alojam nas suas cavidades (\emph{alveolos}) os \emph{dentes}, destinados a mastigar os alimentos.

b) Pharinge. --- A \emph{pharinge} é a parte do canal digestivo, que por um lado confina com a 5\emph{occa}, de que está separada pelo veo \emph{palatino}, e por outro é seguida pelo \emph{esophago}; communiea em cima e por diante com as fossas nasaes, e mais abaixo com a \emph{laringe}, de que está separada pela \emph{epiglotte}. --- A \emph{pharinge} é formada por duas membranas, uma externa de natureza muscular e outra interna mucosa, e é o orgão essencial da deglutição.

c) \textbf{Esophago}. --- O \emph{esophago} é um canal. cylindrico, que se continua com a \emph{pharinge}, desce ao longo da columna vertebral, passa entre os pulmões por traz do coração, atravessa o \emph{diaphragma}, e se abre no \emph{estomago} por um orificio, chamado \emph{cardia}. --- A estructura do \emph{esophago} é identica á da \emph{pharinge}.

d) \textbf{Estomago}. --- \emph{Estomago} é um saeeo membranoso, situado na parte superior do \emph{abdomen}, immediatamente por baixo do \emph{diaphragma}, de um lado continua-se com o \emph{esophago} e por outro com o intestino \emph{delgado}, com o qual communica \pg{13}  por uma abertura, chamada \emph{pyloro}. --- E' formado per tres membranas: una externa \emph{serosa}, outra media \emph{muscular}, e a terceira interna \emph{mucosa}. --- HU o orgão principal da digestão.

e) Intestino delgado. --- O \emph{intestino delgado} é um tubo estreito, que principia no \emph{pyloro}, dobra-se um grande numero de vezes sobre si mesmo e termina no intestino \emph{grosso}. --- Divide-se em tres partes: \emph{duodenum}, que se segue ao estomago, \emph{jejunum} e \emph{ion}, que se seguem ao \emph{duodenum}. --- W' formado de tres tunicas: uma externa \emph{serosa} (\emph{peritoneu}), uma media «\emph{muscular}, e uma interna \emph{mucosa}.

f) Intestino grosso. --- O \emph{intestino grosso} é um canal mais volumoso, que se segue ao intestino \emph{delgado}. --- Divide-se em tres partes: \emph{cecum}, \emph{colom} e \emph{rectum}. --- W" formado, como o intestino \emph{delgado}, de tres membranas \footnote{Ao canal \emph{digestivo} estão \emph{annexos} outros orgãos, com o fim de fabriearem \emph{liquidos} especiaes, necessarios para a digestão. --- Esses orgãos \emph{annevxos} são principalmente as \emph{glandulas salivares}, o \emph{figado} e o \emph{pancreas}.\par \emph{a}) As \emph{glandulas salivares} são formadas por pequenas granulações arredondadas, que têm por fim produzir a \emph{saliva}. São \emph{seis}, dispostas em tres pares: duas chamam-se \emph{sublinguaes}; duas, \emph{submacillares}; e duas, \emph{parctidas}. Estão situadas symetricamente de um e outro lado da bocca, com a qual communicam por uns pequenos canaes excretores.\par \emph{b}) O \emph{figado} é a maior glandula do corpo humano; está situado na parte direita e superior do abdomen, e segrega a \emph{bilis} ou \emph{fel} e a materia \emph{gligogenica}. E' impar e irregular, a sua face superior é convexa, e a inferior é concava. Os canaes secretores da \emph{bilis}, reunidos em ramos, formam o \emph{canal hepatico}, que sai do figado, communiea com a \emph{vesícula} do \emph{fel}, que é o reservatorio da bilis, e vai abrir-se no duodeno. --- A materia \emph{gligogenica} passa do figado para a corrente sanguinea pelo systema da \emph{veia porta}.\par \emph{c}) O \emph{pancreas} é uma glandula situada no abdomen, transversalmente collocada entre o estomago e a columna vertebral. Segrega um liquido, chamado \emph{succo pancreatico}, que vai ao \emph{intestino delgado} pelo \emph{canal} de Vir. \emph{surg}.\par Deseriptos os diversos orgãos que concorrem para a \emph{digestão}, vejamos agora o modo por que esta funcção se realiza.\par \emph{Digestão} é a serie dos actos pelos quaes o homem (e, em geral, o animal) converte os alimentos em materias assimilaveis. Esses actos são (além da \emph{apprehensüo dos alimentos}) \emph{á mastigação}, a \emph{insalivacdo}, a \emph{deglutição}, a \emph{chymificacüo} e a \emph{chylificagdo}.\par O homem \emph{apprehende os alimentos}, in\&rodul-os na bocca; se são liquidos, descem directamente ao estomago; se são solidos, são \emph{mastigados} pelos dentes, ao mesmo tempo que pelas \emph{glandulas salivares} são impregnados de \emph{saliva}. Pela \emph{mastigação} e \emph{insalivagdo}, os alimentos convertem-se em massa molle, chamada \emph{bolo alimentar}. --- Da bocca o \emph{bolo} passa para a \emph{pharinge}, atravessa o \emph{esophago} e chega á cavidade do \emph{estomago}. Chegado alli, transforma-se, pela acção chimica dos succos do estomago, numa pasta semi-liquida e cinzenta, que tem o nome de \emph{c/iymo}. O \emph{chymo} passa do \emph{estomago} para\pgn{14}o \emph{intestino delgado} atravez do \emph{pyloro}, e alli soffre nova transformação pela acção chimica da \emph{bilis} e do \emph{succo:pancreatico}; parte converte-se num liquido branco e leitoso, que se chama \emph{cAylo}, e parte é rejeitada como impropria para a nutrição. Com a formação do \emph{chylo} termina o acto \emph{digestivo}.\par A! \emph{digestão} succede a \emph{absorpção}, que é o acto pelo qual á materia nutritiva, extrahida dos alimentos, passa do canal digestivo para a massa do\par D\pgn{15}sangue, com o fim de regenerar este liquido. --- A \emph{absorpção} faz-se pelas \emph{veias} e pelos \emph{vasos chyliferos}. As \emph{veias} absorvem directamente alguns dos liquidos introduzidos no estomago; --- os \emph{vasos chyliferos}, que vão terminar na membrana mucosa do \emph{intestino delgado}, absorvem as substancias que formam o \emph{cAylo}, \emph{e} que desimboccam no \emph{canal tAoracico}, passando. assim o \emph{chylo} para a veia \emph{sub-clavia esquerda} e entrando na torrente circulatoria.}.

%%>>> t14

\pg{14} \tese{14}{Orgãos do apparelho circulatorio.}
Os orgãos, que constituem o apparelho \emph{circulatorio}, são o \emph{coração}, as \emph{arterias}, OS \emph{vasos capillares} e as \emph{veias}.

a) \textbf{Coração}. --- O \emph{coração} é um orgão muscular, 0co, de forma eoniea, envolvido num saceo fibro-seroso ( '\emph{pericardio}) e situado obliquamente ao lado esquerdo do \emph{thorax}, entre os dois pulmões, com a base para cima e o vertiee para baixo.

FIGURA I. - Coração do homem. \emph{A}, ventriculo direito; --- B, ventriculo esquerdo; --- C, auricula direita; --- D,

auricula esquerda; --- E, arco aortico; --- F, arteria pulmonar; --- G, tronco

brachio-cephalico da aorta; --- H, arterias carotidas primitivas direita e

esquerda; --- \textbar{}, arterias subclavias direita e esquerda; --- K, veia cava -

superior; --- L, veias pulmonares.

No \emph{coração} ha quatro cavidades, --- duas na \emph{base} e são a \emph{auricula direita} e a \emph{esquerda}, --- e duas no \emph{vertice} e são o \emph{ventriculo direito} e o \emph{esquerdo}. --- Quasi no meio ha um septo \pg{15}  vertical, que divide o \emph{coração} em duas metades, cada uma das quaes contém uma \emph{auricula} e um \emph{ventriculo}. A metade, ' que contém a \emph{auricula direita} e o \emph{ventriculo direito}, chama-se \emph{coração direito} ou \emph{pulmonar}; a outra metade, que contém a \emph{auricula esquerda} e o \emph{ventriculo esquerdo}, chama-se \emph{coração esquerdo} ou \emph{aortico}.

As \emph{auriculas} não communicam entre si, \emph{o} que tambem acontece aos \emph{ventriculos}; pois que as duas metades do coração

FIGURA IL --- Córte vertical do coração humano, mostrando o interior

das suas quatro cavidades. \emph{A}, ventriculo direito; --- B, ventriculo esquerdo; --- €, auricula direita; ---

D, auricula esquerda; --- E, orificio auriculo-ventricular direito e valvula

tricuspide; F, orificio aurieulo-ventricular esquerdo e valvula mitral; ---

G, orificio da arteria pulmonar e sua valvula sygmoidea; --- H, orificio

da aorta e sua valvula sygmoidea; --- \textbar{}, orificio da veia cava inferior;

--- K, veia cava superior; --- LL, veias pulmonares.

estão divididas por um septo imperfurado. --- Todavia cada \emph{auricula} communica com o \emph{ocventriculo} do mesmo lado, por meio de um orificio, que se chama \emph{auriculo-centricular}. Cada

um qd estes orificios é fechado por uma valvula, que abaixa quando o sangue passa da auricula para o ventriculo, e depois levanta-se para impedir o refluxo do sangue. A valvula do orifieio \emph{direito} chamar-se \emph{tricuspide}; e a do orifieio \emph{esquerdo} diz-se \emph{mitral}.

As cavidades \emph{direitas} do coração contêm sangue vermelho-eseuro, que se diz \emph{sangue venoso}; as \emph{esquerdas} contem sangue vermelho-rutilante, que se chama \emph{sangue arterial},

b) Arterias. --- As \emph{arterias} são canaes Ou vasos que levam o sangue do coração a todas as partes do corpo.

As principaes arterias são a \emph{pulmonar} e a \emph{aorta}.

\emph{A pulmonar}, que contém Sangue \emph{venoso}, nasce do venirieulo \emph{direito}, sobe para a base do coração, onde se cruza com a \emph{aoría}, dividindo-se em seguida em: dois ramos, um para cada pulmão (\emph{arterias pulmonares} - \emph{direita} e \emph{esquerda}).

\emph{A aorta}, quegeontém sangue \emph{arterial}, nasce do ventriculo \emph{esquerdo}, sobe até á base do coração, curva-se depois para o lado esquerdo, formando o \emph{arco aortico}; em seguida dirige-se para baixo ao longo do thorax, atravessa o diaphragma, segue ao longo do abdomen, aonde termina, dividindo-se em dois ramos, chamados arterias \emph{iliacas primitivas} \footnote{A \emph{aorta}, no seu trajecto, dá origem a outras arterias. Às principaes são :\par a) O \emph{tronco brachio-cephalico}, que depois se divide na \emph{carotida primitiva direita}, e na \emph{sub-clavia direita}; a primeira das quaes leva o sangue para a metade direita da cabeça, e a segunda leva-o para o membro superior direito.\par \emph{b}) \emph{A carotida primitiva esquerda} e á \emph{sub-clavia esquerda}; a primeira das quaes se distribue para a metade esquerda da cabeça, e a segunda para o membro superior esquerdo.\par \emph{cJ} O \emph{tronco ciliaco}, que se divide em tres arterias: a \emph{hepatica} para o figado, a \emph{splenica} para o baco, e a \emph{stomachica} para o estomago. \emph{d}) Às arterias \emph{renaes}, \emph{mesentericas} e \emph{iliacts}, que se distribuem nos rins, nos intestinos e nos membros inferiores.}.

Nos orificios das \emph{arterias}, como tambem nos das \emph{veias}, existem valvulas que impedem o refluxo do sangue, e se chamam \emph{sygmoideas}.

c) Vasos capillares. --- As arterias, afastando-se do coração, dividem-se e subdividem-se em ramos, que são cada vez mais pequenos, e que se enlaçam e anastomosam, formando uma vasta rede, cujas malhas entram em todos os orgãos para lhes levarem o liquido nutritivo. A estas ramificações dá-se o nome de \emph{vasos capillares}. !

Depois de um trajecto mais ou menos longo, os \emph{vasos capillares} reunem-se em canaes cada vez mais amplos e dão \pg{17}  origem ás \emph{veias}: e é d'este modo que as arterias \textbf{communicam} com as veias. --- Nos \emph{vasos capillares} o sangue de vermelho-rutilante torna-se escuro.

d) Velas, --- As \emph{veias} são canaes ou vasos, que levam o sangue de todos os pontos do corpo ao coração

As principaes \emph{veias} reduzem-se a tres systemas:; --- um é constituido pelas \emph{veias cavas}: \emph{superior e inferior}, que, formadas pelas veias de todas as partes do corpo (á excepção das dos pulmões), acompanham a \emph{aorta} e levam ao coração o sangue \emph{venoso}, abrindo-se na \emph{auricula direita}; --- outro é formado pelas \emph{quatro veias pulmonares}, que acompanham as duas arterias pulmonares, e levam ao coração o sangue \emph{arterial}, vegenerado nos pulmões, abrindo-se na \emph{auricula esquerda}; --- o terceiro é o \emph{systema da veia porta}, que é constituido pela reunião das velas \emph{mesentericas}: \emph{superior} e \emph{inferior} e pela veia \emph{splenica}, e que, depois de se ter ramifieado na substancia do figado, se reune e se abre na \emph{cava inferior} \footnote{As \emph{veias}, como tambem as \emph{arterias}, são constituidas por tres tunicas: uma \emph{externa} ou \emph{cellulosa}, muito resistente; outra \emph{media}, formada por fibras elastioas; uma terceira \emph{interna}, mais delgada e formada por um epithelio de cellulas alongadas. A tunica \emph{interna} das veias apresenta um \textbf{grande} numero de prégas, quo servem de \emph{valvulas}, dispostas aos pares e tendo a sua cavidade voltada para o coração, impedindo assim o refluxo do sangue.\par Digamos alguma coisa ácerca do mechanismo da \emph{circulação}.\par a) Circuração. --- \emph{Circulação} é o movimento continuo pelo qual o sangun percorre todas as partes do organismo, irradiando do coração até aos capillares, e voltando depois d'hai novamente ao coração.\par O \emph{ventriculo esquerdo}, contrahindo-se, lança o sangue \emph{arterial} na \emph{aorta}, e por isso em todas as arterias em que ella se divide. O sangue \textbf{passa} depois aos \emph{vasos capillares}, transformando-se alli de arterial em venoso, e d' estes vasos ás veias, que finalmente reunidas formam dois grossos troncos, --- \& \emph{veia cava superior} e a \emph{inferior}, que se abrem na \emph{auricula direita}. Da \emph{auricula direita} o sangue passa ao \emph{ventriculo direito}. O \emph{ventriculo direito}, contrahindo-se, lança o sangue na \emph{arteria pulmonar}, d' onde elle passa aos \emph{pulmões} para se regenerar, e, depois de ter percorrido os capillares d'estes orgãos, volta pelas velas \emph{pulmonares} ao coração, abrindo-se na \emph{auricula esquerda}, d'onde passa ao \emph{ventriculo esquerdo}, para depois recomecar o mesmo trajecto.\par Este trajecto, pois, descreve dois circulos: um, que vai do \emph{ventriculo esquerdo á auricula direita}, e que se chama \emph{grande circulação}; --- outro, que vai do \emph{ventriculo direito á auricula esquerda}, e que se chama \emph{pequena circulação}. --- Na \emph{grande circulação}, as \emph{arterias} conduzem \textbf{sangue} \emph{arterial}, e as \emph{veas} sangue \emph{venoso}; --- nu \emph{pequena circulação}, dà-se o contrario. --- Nas \emph{arterias}, a direcção do sangue é \emph{centrifuga}; nas \emph{veias}, \emph{é centripeta}.\par b) \textsc{CIRCULAÇÃO} \emph{wo} CORAÇÃO. --- O coração, por ser um orgão muscular, tem a propriedade de se contrahir, para em seguida se dilatar. A contracção\par 2 : Vor. II.\pgn{18}chama-se \emph{systole}; a dilatação, \emph{díastole}. O coração não se contrahe nem se dilata todo ao mesmo tempo; mas, quando se contrahem.as \emph{auriculas}, dilatam-se os \emph{ventriculos}; e vice-versa.\par As \emph{auriculas}, estando cheias de sangue trazido pelas \emph{veias cavas} e \emph{pulmonares}, contrahem-se ao mesmo tempo. Em virtude d'esta contracção, o maior volume do sangue é expellido para os \emph{ventrículos}, e só uma pequena quantitade reflue para as veias. Os \emph{ventriculos} dilatam-se para receber o sangue, e, estando cheios, contrahem-se; e o sangue, não podendo refluir para as auriculas pela opposição que lhe fazem as valvulas \emph{auriculo-ventriculares}, penetra nas arterias. Ao mesmo tempo que os \emph{ventrículos} se conirahem, as \emph{awriculas} tornam \& dilatar-se, enchendo-se de sangue, que trazem as veias. Estes movimentos são repetidos e frequentes, e variam com a idade, sexo, alimentação, etc.\par 0) \textsc{CIRCULAÇÃO} NAS ARTERIAS. --- O sangue, nas arterias, caminha do coração para os capillares. --- As causas d'este movimento são as contracções do coração e a elasticidade das arterias. Na verdade, pelas contraeções do coração, o sangue é impellido para as arterias. Mas não basta. As coniracções do coração sao intermittentes, e o movimento é continuo. E' necessaria, pois, uma outra causa. Esta causa é a elasticidade das arterias. Por serem elasticas, as arterias dilatam-se a cada contracção do coração; cessando aquella contracção, contrahem-se, e assim impellem para deante o sangue, que não póde voltar para o coração por causa das valvulas sygmoideas. D'este modo, começando a pressão das arterias sobre o sangue, quando acaba a pressão que sobre elle exerce o coração, segue-se que o movimento do sangue deve ser continuo. --- Podemos, pois, dizer que as arterias estão em \emph{diastole}, quando o coração está em \emph{systole}; e vice-versa. O movimento rythmico da contracção e dilatação do coração e das arterias constitue o phenomeno conhecido do \emph{pulso}.\par d) \textsc{CIRCULAÇÃO} NOS VASOS CAPILLARES. --- O movimento do sangue nos vasos capillares é uniforme e mais lento do que nas arterias. Esse movimento é devido ao impulso que o sangue traz das contracções do coração e das arterias, e á contractilidade propria d'estes vasos.\par e) \textsc{CIRCULAÇÃO} NAS VEIAS. --- O sangue, nas veias, caminha das diversas partes do corpo para o coração. Esse movimento, que é uniforme, é devido não só aos movimentos do coração e das arterias, mas tambem á inspiração, ás contracções musculares, e sobre tudo ás innumeras valvulas que impedem o refluxo do sangue.}.

%%>>> t15

\pg{16} \pg{18} \tese{15}{Orgãos do apparelho respiratorio.}
Os orgãos, que constituem o apparelho \emph{respiratorio}, são a \emph{larynge}, a \emph{trachéa}, os \emph{bronchios}, os \emph{pulmões}, e o \emph{thorazc}.

a) \textbf{Larynge}. --- A /\emph{arynge} é uma cavidade de paredes eartilago-membranosas, aberta superiormente na \emph{pAarynge} e ecommunieando inferiormente com a \emph{trachêa}.

b) \textbf{Trachéa}. --- A \emph{trachêa}, ou \emph{trachéa-arteria}, é um: canal eylindrieo, formado de anneis cartilagineos, ligados entre si por laminas de tecido \emph{conjunctivo} e fibras musculares. Continua-se superiormente com a \emph{larynge}, passa por deante do \pg{19}  \emph{esophago}, até á 4º vertebra dorsal, aonde se divide em dois ramos, que são os \emph{bronchios}.

c) Bronchios. --- Os \emph{bronchios} são os dois ramos, em que inferiormente termina a \emph{trachéa}, e que, entrando nos pulmões, se dividem. successivamente em ramificações cada vez mais delegadas, até terminar em pequenas cavidades, que são as \emph{vesiculas}- \emph{bronchicas}. i

d) Pulmoóes. --- Os \emph{pulmões} são dois orgãos que, junetamente com o coração e os grandes vasos, oceupam a cavidade thoraeiea. --- Têm forma irregularmente conica, cuja base, um

FIGURA III --- Coração e apparelho respiratorio do homem. AA, pulmões direito e esquerdo; --- B, coração envolvido pelo pericardio; ---

C, origem do arco aortico; --- D, arteria pulmonar; --- E, veia cava su-

perior; --- F, trachéa-arteria; --- 68, bronchios; --- HH, \textbar{}l, veias jugula-

res; --- KK, veias sub-clavias; --- P, parte inferior do esterno; --- Q, veia

coronaria anterior; --- R, auricula direita.

pouco obliqua, assenta sobre o \emph{diaphragma}; o seu aspecto é marmoreo, de côr cinzenta rosada; são formados pelos bronchios e pelas \emph{vesiculas bronchicas}, em cujas paredes estão as ultimas ramificações da arteria pulmonar e as primeiras das velas pulmonares. --- A sua superficie é lisa; a do pulmão \emph{direito} apresenta anteriormente dois suleos, que o \textbf{dividem} \pg{20}  em tres /\emph{óbos}; á do pulmão \emph{esquerdo} tem apenas um sulco; que o divide em dois /\emph{óbos}. --- Cada um d'elles está envolvido numa membrana serosa, chamada \emph{pleura}.

€) \textbf{Thorax}. --- O \emph{fAorar}, como dissemos, é uma cavidade que aloja os \emph{pulmões}, o \emph{coração}, e outros orgãos. A sua forma é de uma pyramide conica, tendo o vertice para cima e a base para baixo. É separado do abdomen pelo \emph{diaphragma} \footnote{Deseriptos os orgãos de apparelho \emph{respiratorio}, digamos alguna coisa Acerca da \emph{respiração}.\par a) \textsc{ResPIRAÇÃO}. --- \emph{Respiração} é a funcção que tem por fim converter o sangue \emph{venoso} em sangue \emph{arterial}, pela acção do ar. --- Nesta funcção distinguem-se duas classes de phenomenos: \emph{meclianicos} e \emph{chimicos}.\par b) PHENOMENOS \emph{MECHANICOS} DA \emph{RESPIRAÇÃO}. --- Os phenomenos \emph{mechanicos da respiração} têm por fim determinar a entrada e sahida alternada do ar nos pulmões. Consistem, pois, em dois actos oppostos, que são a \emph{inspiração} e a \emph{expiração}. --- Pela \emph{inspiração}, á cavidade thoracica augmenta; o pulmão dilata-se; o ar interior, augmentando de volume, diminue de tensão; e o ar exterior, vencendo a tensão do ar interior, entra pelas \emph{fossas nasaes} ou pela \emph{bocca}, e precipita-se nas \emph{vesículas Dronchicas}. --- Pela \emph{expiração}, á cavidade thoracica diminue; o ar interior, por diminuir de volume, augmenta em tensão, e, vencendo a tensão do ar exterior, sahe para fóra do pulmão.\par 0) PHENOMENOS CHIMICOS DA \emph{RESPIRAÇÃO}. --- Os phenomenos \emph{chimicos da respiração} consistem nas modificações, que soffrem o \emph{ar} e o \emph{sangue}, postos em contacto nos pulmões.\par a) \emph{Modificações do ar}. O ar \emph{inspirado} compõe-se essencialmente de 20,8 partes de oxygenio para 19,2 de azote. O ar \emph{ewpirado} contém apenas 16,08 de oxygenio. Como se vê, a respiração faz perder ao ar, em cada inspiração, por 100 volumes, 4,77 de oxygenio. Mas, se o ar perde em oxygenio, ganha em anhydrido ecarbonico, porque o ar \emph{expirado} traz, por 100 volumes, 4,26 d'este gaz.\par \emph{b}) \emph{Modificações do sangue}. O sangue, chegando ao coração pela arteria pulmonar, de venoso torna-se arterial. Esta transformação é devida a uma troca gazosa, que se realiza atravez da vesicnlas pulmonares, e pela qual \textbf{entra} o oxygenio para o sangue, e sahe d'elle anhydrido carbonico com o \textbf{vapor} d'agua e algum azote. Na verdade, a quantidade de oxygenio no \textbf{sangue} \emph{arterial} (20 por 100) 6 maior do que no sangue venoso (12 por 100);\pgn{21}e a quantidade do anhydrido carbonico no sangue arterial (84 por 100) é menor que no sangue venoso (48 por 100).\par Além d' esta respiração, que se realiza nos orgãos do apparelho respiratorio, ha no homem uma respiração sypplementar pela \emph{pelle}; pela qual sahem o anhydrido carbonico e o vapor d'agua, e entra o oxygenio.}.

%%>>> art-3

\section*{Artigo III.}
\addcontentsline{toc}{section}{Artigo III.}

\textbf{Orgãos da sensibilidade externa},

%%>>> t16

\tese{16}{Orgãos da sensibilidade externa.}
Os orgãos da \emph{sensibilidade externa} são os que constituem os apparellos dos cinco sentidos \emph{externos}.

%%>>> t17

\pg{21} \tese{17}{Orgãos do apparelho visual.}
Os orgãos, que constituem “o apparelho \emph{visual}, dividem-se em \emph{essenciaes} e \emph{accessorios}. --- Os \emph{essenciaes} são o \emph{globo ocular} e o \emph{nervo optico}; --- OS \emph{accessorios} são as \emph{orbitas}, as \emph{palpebras} e os \emph{musculos oculares}.

FIGURA IV. --- Córte vertical do olho A, cornea trasparente; --- B, humor aquoso; --- €, pupilla;--- D, iris; --- E, crystalino;

--- F, pregas ciliares; --- G, canal de Petit; --- H, esclerotica; --- 1, choroidéa,

--- K, retina; --- L, humor vitreo; --- M, nervo optico; --- N, 0, musculos

rectos, inferior e superior; --- P, museulo elevador da palpebra superior;

--- Q, palpebra superior; --- R, palpebra inferior.

a) Globo ocular. --- O \emph{globo ocular} é um orgão de forma espheroidal, formado por \emph{membranas} e por \emph{meios transparentes}, atravez dos quaes a luz se refraeta.

\emph{a}) As \emph{membranas} são a \emph{esclerotica}, \emph{á cornea transparente}, a \emph{choroidéa} e a \emph{retina}. --- A \emph{esclerotica}, que se chama tambem \emph{cornea opaca}, é a membrana mais externa do \emph{globo ocular}, á que dá a forma; é d'estructura fibrosa, muito resistente e opaca; continua-se anteriormente com a \emph{cornea transparente}, que é uma membrana cireular, de espessura consideravel, e formada por camadas fibrosas. --- Dentro da \emph{esc/erotica} está a \pg{22}  \emph{choroidéa}, de natureza vascular, que é coberta na face interna por uma substancia pigmentar negra, destinada a absorver os raios luminosos, inuteis para a visão, e que no seu bordo anterior forma as \emph{pregas ciliares}. --- Sob a superficie interna da \emph{choroidéa} está a \emph{retina}, substancia molle e esbranquiçada, formada pela expansão do nervo optico, e destinada a receber as impressões da luz.

6) Os \emph{meios transparentes} são o \emph{humor aquoso}, o \emph{crystalino} e o \emph{humor vitreo}. --- O \emph{humor aquoso} é um liquido transparente, formado de agua com pequena quantitade de albumina e alguns saes em dissolução, e situado entre a face posterior da \emph{cornea transparente} e a face anterior do \emph{crystalino}. No meio do espaço, oceupado pelo \emph{humor aquoso}, ha uma especie de diaphragma circular, chamado \emph{iris}, em eujo centro existe uma abertura, chamada \emph{pupilla}. O espaço comprehendido entre a \emph{cornea} e a \emph{iris} chama-se \emph{camara anterior}; e O que fica entre a \emph{iris} e o \emph{crystalino} diz-se \emph{camara posterior}. --- O \emph{erystalino} é uma lente biconvexa transparente, incolor, formado por camadas coneentrieas, cuja densidade e dureza vai crescendo da cireumfereneia para o centro. Está envolvido por uma membrana transparente, chamada \emph{capsula do crystalino}, e fica situado entre a \emph{iris} e o \emph{humor vitreo}. --- O \emph{humor vitreo} é um liquido gelatinoso e diaphano, envolvido numa membrana fina e transparente, chamada \emph{Ayaloidéa}.

b) Nervo optico. --- O \emph{nervo optico}, segundo par dos nervos craneanos, nasce dos hemisperios cerebraes, e depois de atravessar a \emph{esclerotica} e a \emph{choroidéa}, forma com a sua expansão \emph{a retina}.

C) Orbitas. --- As \emph{orbitas} são duas cavidades osseas, que existem na face e alojam os globos oculares.

d) Palpebras, --- As \emph{palpebras} são dois veus membranosos, que cobrem, quasi completamente, o globo ocular na sua face anterior. Sáo formadas por uma cartilagem, coberta externamente pela pelle e internamente per uma membrana mucosa.

e) Musculos oculares, --- Os \emph{musculos oculares} são os que servem para mover as \emph{palpebras} e o \emph{globo ocular}. São \emph{seis}: quatro \emph{rectos} e dois \emph{obliquos} \footnote{\textsc{MxcHaNISMO} DA VISÃO. --- O globo ocular é muito similhante ao instrumento optico, chamado \emph{camara} escura, A \emph{pupilla} é a abertura, pela qual entram os raios luminosos; os «\emph{eios refrangentes} representam a lente que produz a imagem, a \emph{retina} é q plano em que á imagem se desenha.\par A figura V representa .a \emph{theoria da visão}. Do objecto / partem muitos\pgn{23}raios luminosos que, depois de terem entrado pela \emph{pupilla} e atravessado os \emph{meios refrangentes}, produzem na retina \emph{d} a imagem / do objecto /. Esta imagem, porém, é mais pequena que o objecto, e acha-se invertida. A inversão da imagem ns retina é uma consequencia do trajecto dos raios lu-\par FIGURA V. --- Mechanismo da visão\par a, esclerotica; --- 5, choroidéa; --- c, nervo optico; --- \emph{d}, retina; ---' e, con-\par junetiva; --- \emph{f}, cornea; --- 9, centro optico do olho; --- /, humor vitreo;\par --- \emph{i}, erystalino; --- \emph{A}, palpebra inferior; --- //' objecto e sua imagem.\par minosos pelo meio refrangente. --- Mas, comquanto a imagem dos objectos externos se ache invertida na retina, todavia vemos os objectos na sua posição real; e a razão d'isso é porque vemos o objecto não na imagem da retina, mas directamente em si mesmo por meio dos raios luminosos,}.

%%>>> t18

\pg{23} \tese{18}{Orgãos do apparelho auditivo.}
O apparelho \emph{auditivo} \textbf{divide}-se em tres partes: o \emph{ouvido externo}, o \emph{ouvido medio}, e o \emph{ouvido interno}. --- Cada uma d'estas partes consta de varios orgãos.

a) Ouvido externo. --- O ouvido \emph{externo} compõe-se de duas partes: o \emph{pavilhão} ou \emph{orelha} propriamente dicta, e o \emph{canal auditivo externo}. --- O \emph{pavilhão}, que é formado por cartilagens e coberto pela pelle, tem a forma de corneta acustica, para assim poder colligir e concentrar os sons. --- O \emph{canal auditivo externo} é uma especie de tubo, que do pavilhão se dirige para o ouvido \emph{medio}. É parte cartilagineo, parte osseo, e é coberto pela pelle.

b) Ouvido medio. --- O \emph{ouvido medio} é formado pela \emph{caixa do tympano}, cavidade de forma irregular, situada na espessura do \emph{rochedo}, communieando inferiormente com a parte posterior das fossas nasaes por meio de um canal, chamado \emph{trompa d}! \emph{Eustachio}.

O ouvido \emph{medio} está separado do \emph{canal auditivo} por uma membrana tensa, chamada \emph{membrana do tympano}, e communica com o ouvido \emph{interno} por dois orificios, tapados por septos membranosos, e chamados \emph{janella oval} e \emph{janella redonda}. --- Entre o \emph{tympano} e a \emph{janella oval} existem os \emph{ossinhos do ouvido}, que são o \emph{martello}, \emph{á bigorna}, o \emph{osso lenticular} e o \emph{estribo} \footnote{Os \emph{ossinhos do ouvido} chamam-se \emph{martello}, \emph{bigorna}, osso \emph{lenticular} e \emph{estribo}, porque têm similhanca com os objectos, que estes nomes designam.}

c) Ouvido interno. --- O ouvido \emph{interno}, chamado tambem \emph{labyrintho}, e alojado na espessura do \emph{rochedo}, compõe-se de tres cavidades, que são o \emph{vestibulo}, \emph{os canaes semi-circulares}, e o

FIGURA VI. --- Córte do ouvido.

\emph{A}, pavilhão, --- B, canal auditivo externo; --- 6, membrana do tympano; ---

D, caixa do tympano; --- E, bigorna; --- 6, canaes semi-eireulares; ---

H, caracol; --- 1, trompa d' Eustachio; --- M, martelo.

\emph{caracol}. --- O \emph{vestibulo} é a parte central do ouvido \emph{interno}, e communiea pela \emph{janella oval} com o ouvido \emph{medio}, e pela \emph{janella redonda} com o \emph{caracol}. --- Os \emph{canaes semi-circulares} são \textbf{tres} pequenos tubos osseos, situados na parte superior e posterior do vestibulo, no qual se abrem. --- O \emph{caracol} é uma cavidade enrolada em espiral, situada por.deante e por baixo do vestibulo, e dividida interiormente por uma lamina em \pg{25}  duas secções, que são a \emph{rampa externa}, que se abre no vestibulo, e a \emph{rampa interna}, que termina na \emph{joenella redonda}. --- Na lamina, que divide o caracol em duas partes, ramifiea-se o nervo \emph{auditivo} \footnote{A parte essencial do apparelho auditivo é o ouvido \emph{interno}. O ouvido \emph{medio} e o \emph{externo} são partes accessorias, e contribuem para tornar mais perfeitas as impressões sonoras.\par \textsc{MEcHANISMO} DA AUDIÇÃO. --- O mechanismo da \emph{audição} póde facilmente deduzir-se da disposição anatomica do apparelho \emph{auditivo}. --- As vibrações produzidas pelos corpos sonoros são trausmittidas pelo meio ambiente. ao pavilhão do ouvido. Este recolhe-as e dirige-as pelo canal auditivo até á membrana do tympano, que por esses modo entra em vibração. Em seguida, essas vibrações são transmittidas, pelo ar contido no ouvido medio e pela cadeia dos ossinhos, és membranas da janella oval e da janella redonda, que tambem vibram; d'essas membranas communicam-se ao liquido, que enche o ouvido interno, e finalmente aos filetes nervosos do nervo acustico, onde produzem a impressão, que depois é transmittida ao cerebro.}.

%%>>> t19

\pg{24} \tese{19}{Orgãos do apparelho olfactivo.}
O apparelho \emph{o/factivo} resume-se numa membrana mucosa, chamada \emph{pituitaria}, que guarnece internamente as \emph{fossas nasaes}, e que recebe um nervo de sensibilidade especial, --- o \emph{nervo olfactivo}.

As \emph{fossas} nasaes são duas cavidades osseas, abertas na face e separadas uma da outra por um septo mediano vertical. Communicam com o ar exterior por aberturas chamadas \emph{narinas}, com a pharinge na parte posterior, e com as cavidades que existem na espessura dos ossos maxillares superiores, frontal e esphenoide.

A \emph{pituitaria}, que guarnece as \emph{fossas nasaes}, é uma membrana vaseular e humedecida, e apresenta na superficie um grande numero de saliencias, que lhe dão um aspecto avelludado.

O nervo \emph{o/factivo} nasce do cerebro por tres raizes, e, dilatando-se, forma o \emph{bolbo olfactivo}, d' onde partem numerosos filetes, que se distribuem na parte superior da \emph{pituitaria} \footnote{\textsc{MEcHaANISMO} DO OLFACTO. --- Todos os corpos odoriferos espalham no ar uma grande multidão de particulas. As particulas odoriferas penetram com o ar atmospherico nas fossas nasaes em cada inspiração, e, sendo retidas e fixadas pelo muco nasal, impressionam a \emph{pitwitaria}, e produzem á sensação do \emph{olfacto}. :}.

%%>>> t20

\tese{20}{Orgãos do apparelho gustativo.}
O orgão principal do apparelho \emph{gustativo} é a \emph{lingua}. :

A \emph{lingua} é um orgão carnoso e muito movel, constituido por fibras musculares, entrelaçadas em differentes direcções, livre na parte anterior, e ligada na base ao osso hyoide e á \pg{26}  maxilla inferior. --- É coberta de uma membrana mucosa vascular, que apresenta grande numero de proeminencias, chamadas \emph{papillas linguaes}.

O nervo \emph{gustativo} parece ser o \emph{glosso-pharingeo}, cujos filetes se introduzem na \emph{lingua} \footnote{\textsc{MeECHANISMO} DO GOSTO --- As substancias sapidas, dissolvidas em agua ou' saliva, impressionam o nervo \emph{glosso-pharingeo}, cujos filetes se introduzem na lingua e assim produzem a sensação do gosto.}.

%%>>> t21

\tese{21}{Orgãos do apparelho do tacto.}
A séde principal do \emph{tacto} é na superficie palmar das extremidades dos dedos, na parte anterior da bocea e na ponta da lingua \footnote{Não deve confundir-se o \emph{sentido do tacto} com a \emph{sensibilidade tactil}. --- A. sensibilidade \emph{tactil} é mais geral e encontra-se mais ou menos em toda a superficie da pelle e' da mucosa, ae passo que o sentido do \emph{tacto} se acha localisado em certas regiões d' estas membranas.\par Pela sensibilidade \emph{tactil} apenas temos conhecimento das impressões; mas tal conhecimento não é sufficiente para se apreciarem os corpos, que produziram essas impressões. Pelo \emph{sentido do tacto} percebemos não só as impressões, mas tambem as propriedades dos corpos que as produziram, como a forma, o volume, a dureza, o estado da superficie e a temperatura.\par Todavia os nervos da sensibilidade :\emph{actil} são os mesmos que os do sentido do tacto; com a differenca de que as terminações dos nervos tacteis são mais numerosas no sentido do \emph{tacto} do que na sensibilidade \emph{tactil}.\par \textsc{MzcHANISMO} DO TACTO. --- O tacto exerce-se por meio das impressões recebidas nas papillas da \emph{derme}, por intermedio da \emph{epiderme}, ou nas ramifieações dos nervos dà sensibilidade, que se distribuem na superficie d'algumas membranas mucosas e na espessura dos orgãos sensiveis. Estas impressões são depois transmittidas ao cerebro.}.

%%>>> art-4

\section*{Artigo IV.\\ Orgãos da sensibilidade interna e da locomoção.}
\addcontentsline{toc}{section}{Artigo IV. --- Orgãos da sensibilidade interna e da locomoção.}

%%>>> t22

\tese{22}{Orgãos da sensibilidade interna.}
Os orgãos da sensibilidade \emph{interna} são: o \emph{encephalo}, \emph{á medulla espinal} e os \emph{nervos}.

a) Encephalo. --- O \emph{encephalo}, assim chamado por se achar encerrado na caixa craneana, divide-se em tres secções, que são: o \emph{cerebro}, o \emph{cerebello}, e o \emph{bolbo rachidiano} , e que estão ligados entre si pelo \emph{isthmo do encephalo}.

\emph{a}) \emph{O cerebro} é uma massa polposa e molle, formada de \emph{substancia nervosa branca} e \emph{cinzenta}, que occupam - a \emph{branca} Oo \emph{interior} e a \emph{cinzenta} o \emph{exterior} do cerebro, encerrada na \pg{27}  região anterior e na parte posterior da cavidade do craneo, e envolvida em tres membranas, chamadas \emph{menénges} \footnote{A substancia \emph{nervosa branca} é formada principalmente de fibras; occupa quasi todo o interior do cerebro. A substancia \emph{cinzenta} é formada principalmente de cellulas; constitue a camada cortical do cerebro, como tambem a do cerebello; encontra-se em diversos pontos do interior do cerebro.\par A mais externa das \emph{meninges} é de natureza fibrosa, e chama-se \emph{dura-mater}; a media, de natureza serosa, é a \emph{arachnoidéa}; e a interna, de natureza cellular, é a \emph{pia-mater}.}. A suns face superior tem a forma d'um ovoide deprimido.

O cerebro apresenta na sua superficie um grande numero de salieneias ou \emph{circumvoluções}, separadas por sulcos sinuosos ou \emph{anfractuosidades}; e divide-se em duas partes eguaes e \emph{sy}-

FIGURA VII. --- Cerebro humano na sua face superior. \emph{A}, grande \emph{divisão} inter-hemispherica; --- B, sulco de Rolando; --- €, circum-

volugüo frontal ascendente; --- D, cireumvolução parietal ascendente; ---

E, primeira cireumvolucao frontal; --- F, segunda cireumvolucao frontal;

--- 6, terceira circumvolução frontal; --- H, cerebello; --- 1, circumvolu-

ções do lóbo occipital. metrieas, que são os \emph{Ahemispherios cerebraes}. --- Divididos por um sulco profundo, os dois \emph{Hemispherios} unem-se na base por uma fita de substancia nervosa e branca, --- o \emph{corpo calloso}.

Em pontos symetrieos dos dois hemispherios e sob o o \emph{corpo calloso} encontram-se duas especies de ganglios, que são os \emph{corpos estreados} e os \emph{thalamos opticos}. --- Cada um dos \pg{28}  \emph{hemispherios} divide-se em tres \emph{lóbos}: lóbo \emph{anterior} ou \emph{frontal}, lóbo \emph{posterior} ou \emph{occipital}, e lóbo \emph{parietal}, entre o frontal e o oeeipital. --- O lóbo \emph{parietal} toma, na \emph{base}, o nome de \emph{temporal} ou \emph{temporo-esphenoidal} \footnote{O \emph{cerebro} é a séde das mais elevadas faculdades sensitivas. A este orgão convergem, como a centro commum, as impressões sensiveis; e d'elle partem as ordens, que determinam os movimentos voluntarios.\par Na opinião do Dr. Luys, os nervos \emph{sensitivos} de todos os pontos da peripheria reunem-se, sem se confundirem, nos //\emph{i\&lamos opticos}, e d' ahi se espalham pars todos os pontos da camada cortical, para alli levar e conservar as sensações; --- os nervos \emph{motores} de todos os pontos da camada cortical reunem-se nos \emph{corpos estreados}, para levar alli os movimentos volontarios, que depois devem communiecar-se aos musculos.\par O emispherio \emph{direito} produz os movimentos da parte \emph{esquerda} do corpo; e vice-versa. E é por isso que, em certos casos pathologicos, a parte doente do cerebro é ordinariamente opposta á parte paralysada do corpo. Todavia não é sempre assim; porque está averiguado que um só hemispherio basta para determinar o movimento de todo o corpo.\par Os dois hemispherios cerebraes funccionam symetricamente, estão numa reciproca dependencia, um parece a repetição de outro. É Por isso que no homem a destruição \emph{lenta} e \emph{progressiva} de um hemispherio passa desapercebida durante a vida, exercendo o hemispherio são as funcções do hemispherio doente.\par Alguns physiologistas modernos opinaram que o cerebro é absolutamente indispensavel para \emph{todas} as manifestações da vida sensitiva. Mas não é assim, Be se tirarem os hemispherios cerebraes a um animal, este mostra que ainda vê, ouve, sente... , embora não saiba interpretar as suas sensações; perdeu a memoria, a imaginação, a espontaneidade, o instincto, e, numa palavra, todas as faculdades superiores sensitivas. Por isso os hemispherios cerebraes centralisam e aperfeicóüam as sensações, mas comtudo não são o seu orgão immediato e essencial,}.

5) O \emph{cerebello} é uma massa menos volumosa do que o cerebro, tem forma elliptica, coberto pela parte posterior do cerebro, de que está separado por um prolongamento das meninges, chamado \emph{pequena fóice do cerebro}. É formado de tres lóbos, dois lateraes e symetricos, e um terceiro mediano entre os lateraes.

O cerebello continua-se com a \emph{medulla espinal} por tres especies de raizes, chamadas \emph{pedunculos do cerebello}. As raizes \emph{medias} unem os dois lóbos do cerebello e tomam o nome de \emph{ponte de Varolio}.

O cerebello, como o cerebro, é constituido de \emph{substancia branca} e de \emph{substancia cinzenta}; mas estas duas substancias estão dispostas por tal modo que, fazendo um eórte transversal \pg{29} , a superficie da secção apresenta um aspecto arboreo, a que se dá o nome de \emph{arvore da vida} \footnote{As funcções physiologicas do \emph{cerebello} são ainda bastante escuras. \emph{Milne-Edwards} diz que o cerebello é necessario para a produecüo harmonica ou pira o equilibrio dos movimentos, pois que a experiencia mostra que a lesão d'este orgão importa um desarranjo consideravel na locomoção. Outros contestam as affirmações de Milne-Edwards, e apresentam casos de individuos, que, embora privados do cerebello, caminhavam regularmente.}.:

c) O \emph{bolbo rachidiano} é uma porção da massa encephaliea, immediatamente situada por baixo do \emph{cerebello}, a qual superiormente se une aos dois hemispherios cerebraes por dois peduneulos, chamados \emph{pedunculos cerebraes}, e inferiormente se continua com a \emph{medulla espinal}, de que é um prolongamento \footnote{O \emph{bolbo rachid:ano} exerce as mesmus funcções que são proprias da \emph{medulla espinal}, da qual é uma continuação. Além d'estas funcções, o \emph{bolbo} tem uma grande importancia na producção de certos movimentos, necessarios para a vida nutritiva. Se se picar com um alfinete um pequeno sulco em forma de V, que se encontra no punto em que a medulla espinal penetra no craneo, cessam de repente os movimentos do thorax, do diaphragma, e todos os mais indisps»nsaveis para a respiração, e o animal cahe fulminado. \emph{Flourens} dà a este centro o nome de nó \emph{vital}.}.

FIGURA VIII --- Córte vertical do cerebro humano. \emph{A}, hemispherio esquerdo; --- B, corpo calloso; --- €, thalamos opticos; ---

D, protuberaneia annular; --- E, medulla allongada; --- F, córte do cere-

bello mostrando a arvore da vida; --- 6, hemispherio esquerdo do cerebello.

b) Medulla espinal. --- A \emph{medulla espinal} é um longo cordão de substancia nervosa, \emph{cinzenta} e \emph{branca} (sendo esta \emph{externa}), e envolvida pelas tres meninges que se prolongam do cerebro. Continuando-se superiormente com o bolbo \emph{racAidiano}, á \pg{30}  \emph{medulla espinal} occupa o comprimento d'um canal, chamado tambem \emph{racAidiano}, desde a cabeça até á bacia, onde se divide em muitas ramifieacaes, que formam a \emph{cauda equina}. --- No seu trajecto dá origem a numerosos pares de \emph{nervos} \footnote{A. funecüo principal da \emph{medulla espinal} consiste em recolher as impressões dos nervos para as transmittir ao cerebro que as centralisa e con- S?rva, e em íransmittir aos nervos as excitações motrizes e voluntarias do cerebro. O Dr. Béclard diz que a \emph{medulla espinal} é um centro nervoso, como é o cerebro; de modo que, se se cortar num ponto qualquer da sua extensão, a parte separada do cerebro não fica totalmente paralysada.}.

c) Nervos. --- Os \emph{nervos} são cordões de substancia branea, formados de filamentos medullares, que nascem symetrica-

FIGURA IX. --- Cerebro humano na sua face inferior \emph{A}, grande \emph{divisão} inter- hemispherica; --- B, \emph{fenda} de Sylvius; --- €, face

inferior do lóbo frontal; --- D, lóbo esphenoidal; --- E, protuberancia

annular; --- F, bolbo rachidiano; --- 6, lóbo esquerdo do cerebello, ---

H, lóbo medio do cerebello; --- \textbar{}, nervos olfactivos (1.º par.); --- K, nervos

opticos (2º par); --- L, pedunculos cerebraes; --- M, nervos motores ocu-

lares communs (3.º par); --- N, nervos patheticos (4.º parj; --- O, nervos

irigemoos (5.º par); --- P, nervos oculares motores externos (6.º par); ---

Q, nervos faciaes (7.º par) --- R, nervos acusticos (8.º par); --- S, nervos

glosso-pharyngeos (9.º par); --- T, nervos pneumo-gastricos (10.º par);

.--- U, nervos espinaes (11.? par), --- V, nervos grandes hypoglossos (12.9 par). mente d'um e d'outro lado do \emph{cerebro} e da \emph{medulla espinal}, e terminam por divisões extremamente tenues nos orgãos em que se distribuem. Ás vezes reunem-se uns aos outros, formando redes que se chamam \emph{plexos}; nalguns pontos engrossam, e esses engrossamentos denominam-se \emph{ganglios nervosos}.

n --- Uns presidem ás funecóes de \emph{relação}, outros ás da \emph{vida nutritiva}. Os primeiros são os \emph{craneanos} e os \emph{espinaes} ou \emph{rachidianos}; os segundos são os que constituem o \emph{grande sj sympathico}.

\emph{a}) Os nervos \emph{craneanos} nascem na parte inferior do \emph{encephalo}, \emph{e} são 12 \emph{pares}, que tém os seguintes nomes: 1.? par, os nervos \emph{olfactivos}; 2.º, os \emph{opticos}; 3.\textasciicircum{}, os \emph{motores} pde \emph{communs}; 4.º, os \emph{patheticos}; 5.º, os \emph{trigemeos} ou \emph{trifaciaes}; 6.º, 0 \emph{motores oculares externos}; 7.º, \emph{os faciaes}; 8.º, os \emph{acusticos}; 9.º, os \emph{glosso-pharyngeos} ; 10.*, os \emph{pneumo-gastricos}; 11.*, os \emph{espinaes}; 19.?, os \emph{grandes hypoglossos}. D' estes nervos uns são \emph{sensitivos}, e outros são \emph{motores} \footnote{Os nervos \emph{motores} determinam a contracção dos muscolos, em que se distribuem os movimentos voluntarios; os nervos \emph{sensitivos} recebem as impressões sensiveis, e as transmittem ao cerebro.\par São \emph{sensitivos}: os \emph{opticos}, os \emph{alfactivos}, os \emph{acusticos}, etc. ---- São \emph{motores}: os \emph{faciaes}, os \emph{motores oculares}, os \emph{patheticos}. ete.}.

\emph{b}) Os nervos \emph{rachidianos} são em numero de 31 \emph{pares}, e chamam-se \emph{cervicaes}, \emph{dorsaes}, \emph{lombares} e \emph{sagrados}, segundo nascem da região \emph{cervical}, ou da \emph{dorsal}, ou da \emph{lombar}, ou da \emph{sagrada}. Cada um d estes nervos tem fibras \emph{sensitivas} e fibras \emph{motrizes}.

\emph{c}) O \emph{grande sympatico} é um duplo cordão nervoso, formado por duas series de \emph{ganglios}, ligados entre si por filetes nervosos, e situado ao longo. da columna vertebral. Os nervos do \emph{grande sympathico} distribuem-se em todos os orgãos da \emph{nutrição}, no \emph{coração}, nos \emph{pulmões}, nos \emph{orgãos digestivos}, etc.

%%>>> t23

\pg{31} \tese{23}{Orgãos da locomoção.}
Os orgãos da \emph{locomoção} São os \emph{ossos} e os \emph{musculos}. Estes são orgãos \emph{activos}; aquelles, \emph{passivos}.

a) Ossos. --- Os \emph{ossos} são orgãos formados de tecido osseo, vasos sanguineos, nervos e vasos lymphaticos, e protegidos por uma membrana externa, chamada \emph{periosteo}. --- Distinguem-se em \emph{longos}, \emph{chatos} e \emph{curtos}. --- Estão cobertos de partes molles, e reunidos entre si, formando \emph{articulações}.

b) Musculos. --- Os \emph{musculos} são orgãos de côr encarnada, que formam o que vulgarmente se chama \emph{carne}, e que se encontram em todos os pontos em que ha movimento. --- Dividem-se em \emph{lisos} e \emph{estreados}. São estes ultimos os orgãos \emph{activos} da \emph{locomoção} \footnote{\textsc{MxkcHANISMO} DA \emph{LOCOMOÇÃO}. --- Os \emph{musculos} recebem uma excitação da parte dos nervos motores. A esta excitação os musculos, por serem \emph{trritaveis}, respondem com uma \emph{contracção}, pela qual augmentam de espessura, mas diminuem de comprimento. A \emph{contracção} determina o movimento dos ossos e do corpo inteiro. D'este modo os ossos são verdadeiras alavancas submettidas ás leis ordinarias da mechanica (\emph{Langlebert}, \emph{Gervais}).}.

ALMA HUMANA \textsc{SuMMARIO}: --- Alma humana, sux existencia e substancialidade. --- Simpli-

cidade da alma humana. --- Espiritualidade da alma humana. --- A

espiritualidade da alma humaaa e o materialismo.

%%>>> art-5

\section*{Artigo I.}
\addcontentsline{toc}{section}{Artigo I.}

\pg{32} \textbf{Alma humana}, \textbf{sua existencia}

e \textbf{substancialidade}.

%%>>> t24

\tese{24}{Alma humana.}
Alma \emph{humana} é o \emph{principio primeiro pelo qual o homem vegeta}, \emph{sente e entende}. O principio \emph{segundo} das operações é constituido pelas \emph{faculdades}. --- A verdade d'esta, definição, x demonstrada nos artigos seguintes \footnote{A esta parte*da Anthropologia, que trata da \emph{alma Jmana}, dá-se commummente o nome de \emph{Psychologia}. --- Cf. S. Thomaz, \emph{Sum}. \emph{Theol}., I jte q. (5. --- \emph{Contra gent}., l. Il, cc. 68, 64, 65.\par O systema, que se inventou para desvirtuar a \emph{natureza} da alma humana e as suas mais nobres prerogativas, é o \emph{materialismo}. Embora variado e diverso nas suas feições accidentaes e secundarias, segundo a variedade e diversidade dos tempos, dos paizes e dos costumes, este systema conservou sempre, atravez de todas as edades e de todas as mudanças, a mesma essencia, o mesmo fundo. Todos os materialistas concordam em que a alma humana, --- ou não é uma substancia distincta do corpo nem das suas qualidades, --- ou, se é distincta do corpo, é alguma coisa extensa ou dependente da materia nas suas operações e por isso na sua existencia.\par D'ahi a divisão do \emph{materialismo} em \emph{atomico}, \emph{physiologico} e \emph{sensualista}.\par \emph{a}) O materialismo \emph{atomico} sustenta que a alma, embora distincta do nosso corpo, todavia é alguma coisa corporea, composta de atomos e dotada de extensão. --- Este materialism» foi ensinado, na antiguidade, pelos philosophos da Eschola Jonica, por Democrito e Epicuro, e pelos Stoicos.\par 5) O materialismo \emph{physiologico}, ou \emph{dinamico}, ensina que a alma não se distingue do corpo, mas é o temperamento ou a harmonia do proprio corpo, ou uma funcção do cerebro; porque a materia é essencialmente activa e capaz de produzir todas as operações que o homem produz. --- Este materialismo foi sustentado --- na antiguidade, por Dicearcho, Aristoxemo e Galeno, --- no seculo passado, por Holbach, Lamettrie, Cabanis; --- no seculo presente, por Biichner, Molleschott, Wagner, Comte, Vogt, etc.\par \emph{c}) O materialismo \emph{sensualista} admitte que a alma é distincta do corpo, que é uma substancia simples e inextensa, mas nega que ella seja espiritual, isto é, independente da materia. Este materialismo teve por defensores Locke, Condillae, Robinet e os que pertencem á eschola \emph{sensualista}.\par Nos artigos d'este capitulo demonstraremos --- contra o materialismo \emph{Jhysiologico} ou \emph{dynamico} que a alma humana é \emph{distincta} do \emph{corpo}, --- contra o materialismo \emph{atomico} que a alma é \emph{simples} ou \emph{inewtensa}, --- contra o materialismo \emph{sensualista} que ella é \emph{espiritual}.}.

%%>>> t25

\pg{33} \tese{25}{A alma humana é ente real.}
O homem vegeta, sente e entende; e por isso deve possuir um principio primeiro, de que emanam aquellas operações; pois todo o effeito tem uma causa. Esse principio deve ser ente \emph{real}; porque as operações do homem são entes reaes, e todo o effeito exige uma causa proporcionada. Logo a alma humana é ente \emph{real} \footnote{Ninguem nega a \emph{existencia} da alma humana; e ninguem a pode negar sem contradiccio; porque quem nega, pensa; e por isso admitte implicitamente um principio pensante, que é a \emph{alma}. --- Mas, se ninguem nega a \emph{existencia} da alma humana, muitos desvirtuam a sua natureza.}.

%%>>> t26

\tese{26}{A alma humana é substancia.}

a) Todo o ente real ou é \emph{substancia}, ou é \emph{accidente}. Ora a alma humana é o principio \emph{primeiro} da nossa vida, e o principio \emph{primeiro} da vida não póde ser \emph{accidente}. Na verdade, o \emph{accidente}, assim como não póde existir por s; mesmo, mas precisa de um sujeito que o sustente, tambem não póde ope- 'ar por virtude propria, mas por virtude do mesmo sujeito a que pertence, e por isso não póde ser o principio \emph{primeiro} da vida, mas só o \emph{segundo}. Logo a alma \textbf{humÉna} é \emph{substancia}.

b) A alma humana é o principio que constitue o homem numa determinada especie, distinguindo-o de todos os outros seres. Ora este principio não póde ser \emph{accidente}; aliás a differença entre o homem e os outros seres seria apenas accidental, o que é absurdo. Logo á alma humana é \emph{substancia} \footnote{Esta conclusão é contestada pelos \emph{Aantístas}, pelos \emph{positivistas}, e por \emph{Descartes}, \emph{Storchenau} e \emph{Galluppi}.\par a) Os \emph{Aantistas} dizem que as \emph{substancias} (\emph{nowumenos}) escapam á nossa percepção e que só podemos conhecer os \emph{phenomenos}.\par Respondemos que, atravez dos \emph{phenomenos} ou \emph{modificações}, que percebemos \emph{directamente}, deve necessariamente existir uma realidade, que se modifica e a que damos o nome de \emph{substancia}.\par b) Os \emph{sensualistas} affirmam que a alma é um grupo de sensações.\par Esta opinião tambem é falsa. Um grupo de sensações não póde existir sem um sujeito em que resida. De mais, as sensações passam e a alma fica. Logo a alma não é um grupo de sensações.\par 0) Os \emph{positivistas} não reconhecem a \emph{substancialidade} da alma e só admittem \emph{phenomenos}; porque, dizem, pela experiencia só percebemos os phenomenos.\par Tambem esta opinião não póde acceitar-se. Na verdade, pela: experiencia conhecemos \emph{directamente} os phenomenos, mas \emph{indirectamente} a substancia de que derivam; porque conhecemos as sensações como \emph{proprias} e por isso pertencentes a um determinado \emph{sujeito}, e este sujeito é a substancia.\par 4) \emph{Descartes} collocou a essencia da alma no \emph{pensamento}. --- O philosopho francez foi seguido por V. \emph{Gioberti} (\emph{Protol}., Saggio 1, 8 1), que faz consistir a essencia da alma na intuição primitiva de Deus, --- e por A. \emph{Rosmini} (\emph{Nuovo}\par 3 Vor. II.\pgn{34}\emph{Saggio}, sez. V, part. 2, c. 4, a. 1, 8 8), ensinando que \& \emph{intuição do ente possivel é o acto primeiro} que constitue a \emph{intelligencia humana}.\par É falso que a essencia da nossa alma consista no \emph{pensamento}. Na verdade, o pensamento é um \emph{accidente}; ao passo que a \emph{alma}, como foi demonstrado, é \emph{substancia}. --- De mais, o. pensamento deriva da alma. Logo, se a essencia de alma consistisse no pensamento, a causa identificar-se-ia com o effeito, o que é absurdo. --- Finalmente, o pensamento é variavel, multiplo, mutavel;ao passo qua a essencia da alma é constante, una, immutavel. Logo o pensamento não póde constituir a essencia da alma.\par e) \emph{Storchenau} e \emph{Galluppi}, seguidos por alguns esceriptores, collocaram a essencia da alma na \emph{faculdade} de \emph{pensar}.\par Tambem esta opinião é inacceitavel. Com effeito, a faculdade de pensar é um acto \emph{segundo} e \emph{accidental}, e por isso deriva da essencia da alma, que é o acto \emph{primeiro} e \emph{substancial}. Logo a essencia da alma não póde identificar-se com a faculdade de pensar.}.

%%>>> t27

\pg{34} \tese{27}{A alma humana é substancia incompleta.}
Toda a substancia, que é \emph{parte} de um todo, é \emph{incompleta}. Ora a alma humana é \emph{parte} de um todo, que é o homem; porque o homem não é só alma nem só corpo, mas é o composto de alma e de corpo. Logo a alma humana é substancia \emph{incompleta} \footnote{A alma humana, por ser forma \emph{substancial} do composto humano, é substancia \emph{incompleta}, como o são todas as formas \emph{substanciaes}. "Todavia, como veremos, ha differenca entre a alma humana e as outras formas \emph{substanciaes} inferiores; pois aquella póde exercer \emph{algumas} operações independentemente do corpo e por isso é \emph{subsistente}, ao passo que estas não podem exercer operação alguma sem o corpo e por isso não são \emph{subsistentes}.}.

%%>>> t28

\tese{28}{A alma humana é distincta do corpo.}

a) A alma é o principio \emph{primeiro} da vida, isto é, é o principio que communiea a vida ao corpo, e o torna \emph{vivo}, capaz de exercer actos vitaes. Ora o nosso corpo não pode ser o: principio \emph{primeiro} da vida; porque o corpo é vivo não emquanto é \emph{corpo} (aliás todos os corpos viveriam e seriam principios de vida), mas emquanto é corpo \emph{humano}, e É corpo \emph{humano} pela \emph{alma}, que por isso deve ser uma coisa distincta. Logo a alma humana é distineta do corpo \footnote{Em poucas palavras: O corpo \emph{vivo} differe do corpo \emph{no-vivo}, porque\pgn{35}um tem alma, e outro não a tem. Logo a alma, por ser a razão da differenca entre varios corpos, não é corpo. --- \emph{Of}. \emph{Sum}. \emph{theol}., I p., q. 75, a. 1.}.

b) A alma humana, se não fosse distincta do corpo, devia estar sujeita a todas as vicissitudes e mudanças, a que o mesmo. corpo está sujeito. Ora este consequente é falso. Na verdade, ao passo que o nosso corpo soffre, pelas leis de assimilação e excreção, continuas mudanças, de modo que no fin de alguns annos se renova completamente, a alma, que é o principio da vida, não se muda nem se renova, mas fica sempre a mesma; pois a consciencia attesta que nós, que \pg{35}  vivemos hoje, somos os mesmos que viviamos no passado. Logo a alma humana é distineta do corpo \footnote{O celebre \emph{Flourens} conclue as suas admiraveis experiencias com estas palavras: « A materia toda apparece e desapparece, faz-se e desfaz-se, e uma só coisa permanece: é a que faz e desfaz, é a que produz e destroe, é a força que vive no meio da materia e que a governa. » (\emph{De la} Vie \emph{et de V Intell}. , p. 21).\par Alguns materialistas antigos e modernos negam que a alma seja distincta do corpo, e dizem que ella é apenas o \emph{temperamento} do corpo, ou a \emph{harmonia} que existe nos elementos e orgãos do mesmo corpo.\par \emph{a}) A alma humana não é o \emph{temperamento} do corpo. Na verdade, o \emph{temperamento} resulta principalmente do equilibrio das forças physicas e chimicas. Ora \& alma humana, por ser o principio da vida no homem, é distincta das forças physicas e chimicas da materia. Logo a alma humana não é o \emph{temperamento} do corpo. --- De mais, a alma regula e reprime o \emph{temperamento} e as paixões que d'elle derivam. Ora isto seria impossivel, se a alma consistisse no \emph{temperamento}.\par \emph{b}) A alma humana não é a \emph{harmonia} que existe nos elementos e orgãos do corpo. Com effeito, a \emph{Harmonia} é commum a todos os corpos, mesmo aos que não são dotados de vida, mas a alma não se encontra em todos os corpos; --- a \emph{harmonia} augmenta e diminue, aperfeicóa-se e altera-se, ao passo que a alma é indivisivel e immutavel; --- aquella /\emph{armonia} é uma resultante, incapaz de resistir ás paixões e dominar o corpo, mas a alma é o principio, a que todas as forças do corpo estão sujeitas.}.

%%>>> art-6

\section*{Artigo II.\\ Simplicidade da alma humana.}
\addcontentsline{toc}{section}{Artigo II. --- Simplicidade da alma humana.}

%%>>> t29

\tese{29}{Simplicidade da alma humana.}
\emph{Simplicidade} é á propriedade, pela qual a alma humana não é composta de partes \emph{essenciaes}, isto é, de \emph{materia} e de \emph{forma}, como é o homem, nem de partes \emph{integrantes}, como o é toda a substancia extensa \footnote{\emph{Simplicidade}, em geral, é a ausencia de partes ou de composição. --- A \emph{simplicidade}, por si, não significa perfeição nem imperfeição, mas prescinde de uma e outra coisa; porque ha entes \emph{simples} mais perfeitos do que os \emph{compostos} (e isto dá-se quando o simples denota um \emph{todo}, uma coisa \emph{perfeita}, um ser \emph{em acto}), assim o Anjo é mais perfeito do que o homem, --- e ha entes \emph{compostos} mais perfeitos do que os \emph{simples} (e isto dá-se quando o \emph{simples} exprime uma \emph{parte}, uma coisa \emph{imperfeita}, um ser em \emph{potencia}), assim o composto humano é mais perfeito do que a alma só. --- Divide-se\pgn{36}a \emph{simplicidade} em \emph{absoluta} e \emph{relativa}, segundo exclue toda e qualquer composição, ou apenas alguma especie d'esta. --- A \emph{simplicidade}, que attribuimos á alma humana não é a \emph{absoluta}, porque a nossa alma é composta de essencia e de existencia, de substancia e accidente, etc.; mas é a \emph{relativa}, porque excluimos da alma a composição \emph{essencial} e a \emph{integrante}.\par A. \emph{simplicidade} da alma, embora pareça denotar uma coisa \emph{negativa}, emquanto significa \emph{indivisibilidade}, todavia na \emph{realidade} representa uma perteição \emph{positiva}. Com effeito, a \emph{simplicidade} da alma não é mera negação, como é a simplicidade do ponto mathematico, mas é uma propriedade pela qual a substancia da alma não é divisivel, porque não se compõe de partes \emph{essenciaes}, nem de partes \emph{integrantes}.}.

30. A \textbf{alma} humana é \textbf{essencialmente} simples, \textbf{isto} é, não é \textbf{composta} de \textbf{partes} essenciaes. \textbar{}

a) Se a alma humana fosse composta de partes \emph{essenciaes}, isto é, de \emph{materia} e de \emph{forma}, seria ella uma substancia \pg{36}  \emph{completa} na sua especie, e n\&o poderia ser \emph{parte} do composto humano. Ora a nossa alma não é uma substancia \emph{completa} na especie, mas é apenas uma \emph{parte}, embora a mais nobre, do

74 homem. Logo a alma humana é \emph{essencialmente} simples \footnote{Alem d' isso, se a alma humana fosse composta de duas ou mais partes \emph{essenciaes}, estas partes seriam --- ou todas \emph{activas}, --- ou todas \emph{passivas}, --- ou uma \emph{activa} e outra \emph{passiva}. --- Na primeira hypothese, haveria no homem tantas almas, quantas fossem aquellas partes; o que é contrario ao testemunho da consciencia. --- Na segunda, não haveria alma no homem; porque a alma, sendo o principio primeiro da vida, é essencialmente \emph{activa}. --- Na terceira, só a parte \emph{activa} seria a alma, e assim ficaria demonstrada a proposição que affirma a simplicidade \emph{essencial} da nossa alma. --- Cf S. Thomaz, \emph{Sum}. \emph{Theol}., 1 p., q. 75, a. 5.}.

b) O homem, por meio das suas faculdades \emph{sensitivas}, percebe os objectos materiaes com um acto \emph{unico} e \emph{completo}. Ora tal percepção seria impossivel, se em nós não existisse um principio \emph{essencialmente} simples ou inextenso, que recolhesse e unisse a percepção de cada uma das partes das faculdades sensitivas. Na verdade, sem este principio, uma parte da faculdade, embora podesse perceber uma parte do objecto, não poderia perceber ao mesmo tempo \emph{todo} o objecto; e assim a percepção não seria unica nem completa, mas haveria tantas percepções incompletas quantas as partes da faculdade sensitiva; o que é contrario ao testemunho da consciencia. Ora o principio pelo qual o homem sente é a alma. Logo a alma humana é \emph{essencialmente} simples \footnote{A' mesma conclusão chegamos, se considerarmos que, muitas vezes, temos, ao mesmo tempo, sensações \emph{diversas} e até \emph{oppostas}; assim, ao mesmo tempo, sentimos frio nos pés e calor nos mãos, e todavia é o mesmo sujeito que diz: eu sinto frio nos pés e calor nas mãos. Ora isto não se poderia dizer, se, no homem, não existisse um principio \emph{simples} e \emph{inectenso}, capaz de recolher esses sensações oppostas. Diz S. Thomaz: « \emph{Omne divisibile indiget aliquo continente et uniente partes eius} n (\emph{Contra gent}., \textbar{}. lI, c. 65,\par Poderia objectar-se: Para explicar os phenomenos da sensação, não é necessario admittir em nós um principio simples e inextenso, porque, se\pgn{37}uma parte da faculdade sensitiva só pode perceber uma parte do objecto, \emph{toda} a faculdade poderá: perceber \emph{todo} o objecto.\par Resposta. Concedemos que \emph{toda} a faculdade póle perceber e percebe \emph{todo} o objecto, mas repetimos que esta percepção seria impossivel, se a mesma faculdade não fosse informada por um principio inextenso. Na verdade, embora cada uma das partes da faculdade \emph{sensitiva} corresponda a uma parte do objecto, todavia o homem percebe, com um acto \emph{unico} e \emph{completo}, \emph{todo} o objecto. Ora isto seria impossivel,se, na faculdade sensitiva, não existisse um principio \emph{simples} ou \emph{inewtenso}, que recolhesse e unisse as sensações de \emph{cada uma} das partes da faculdade \emph{sensitiva}.\par Insistem : Toda a faculdade sensitiva é extensa, pois reside num orgão material. Logo o principio da vida sensitiva não é \emph{essencialmente} simples.\par Resposta. A faculdade sensitiva, que resulta composta do orgão e do principio vital, é extensa e material emquanto ao proprio orgão, mas não emquanto ao principio vital, que informa o orgão. Este principio vital deve ser \emph{essencialmente} simples ou inextenso; aliás não poderia a percepção da faculdade sensitiva ser, como realmente é, unica e completa.}.

c) A alma humana entende objectos simples ou inextensos, como a \emph{verdade}, a \emph{virtude}, a \emph{justiça}, etc. Ora, se a nossa alma fosse \emph{essencialmente} composta, não poderia entender aquelles objectos. Por quanto, --- ou cada parte da alma perceberia \emph{todo} o objecto, --- ou uma parte da alma só perceberia uma parte do objecto. A primeira hypothese é falsa; porque admitte no homem muitos principios pensantes. A segunda tambem é falsa; porque um objecto simples não póde ser percebido por partes. Logo a alma humana é \emph{essencialmente} simples \footnote{A. \emph{simplicidade} da alma humana demonstra-se não só pela \emph{percepção} dos objectos \emph{simples}, mas tambem pelo \emph{juizo} e pelo \emph{raciocinio}.\par \emph{a}) O \emph{juizo} consiste na affirmação da \emph{conveniencia} ou da \emph{desconveniencia} do \emph{predicado} com o \emph{sujeito}. Para esta affirmacao é necessario que o principio pensante conheça o sujeito, o predicado e a relação que existe entre um e outro. Mas, se o principio pensante fosse composto de partes \emph{essenciaes}, uma parte perceberia o sujeito, outra o predicado, e uma terceira a relação, mas nunca uma e a mesma parte perceberia o sujeito, o predicado e, muito menos, a relação entre este e aquelle. Logo o principio pensante, que é a alma, deve ser \emph{essencialmente} simples.\par \emph{b}) O \emph{raciocinio} tambem demonstra a \emph{simplicidade} da alma. Na verdade, a alma, quando raciocina, deve perceber as premissas e a conclusão. Ora « esta percepção seria impossivel, se a alma fosse composta de partes \emph{essenciaes}; porque então uma parte perceberia uma premissa, outra parte outra premissa, e outra a conclusão, mas nunca uma parte poderia perceber as premissas e a conclusão. Logo é necessario admittir que a alma é \emph{essencialmente} simples, porque. só assim poderá perceber as premissase a conclusão,}.

%%>>> t31

\pg{37} \tese{31}{A alma humana é extensivamente simples, isto é, não é composta de partes integrantes.}

a) Todo o ente que é dotado de extensão, e por isso de \pg{38}  partes \emph{integrantes}, é corpo. Ora a alma humana não é corpo, mas é distincta do corpo. Logo a alma humana é \emph{extensivamente} simples, isto é, não é composta de partes \emph{integrantes}.

b) A alma humana é \emph{essencialmente} simples, isto é, não é composta de partes \emph{essenciaes}. Ora um ser, que é \emph{essencialmente} simples, é-o tambem \emph{extensivamente}; pois de uma essencia \emph{simples} não póde emanar uma propriedade \emph{composta}, como é a \emph{extensão}. Logo a alma humana é \emph{ezstensivamente} simples \footnote{Provam a \emph{simplicidade} da alma todos os argumentos com que no artigo seguinte demonstraremos a sua \emph{espiritualidade}; visto que todo o ente \emph{espiritual} é \emph{simples}, embora nem todo o ente \emph{simples} seja \emph{espiritual}.\par Objectam: Uma substancia \emph{simples} não pode estar num ser extenso. Logo a alma, se fosse \emph{simples}, não poderia estar no corpo.\par Resposta. Uma substancia \emph{simples}, se não póde estar num corpo por modo \emph{circwmscriptivo}, póde estar nelle por modo \emph{definitivo}, emquanto com a sua força sustenta e rege o proprio corpo.\par Insistem : Se a alma fosse \emph{simples}, não poderia mover o corpo; porque um ser extenso só póde ser movido por um agente extenso.\par Resposta. Ha um duplice contacto: de \emph{força} e de \emph{quantidade}. Uma substancia \emph{simples}, como é á nossa alma, não póde mover o corpo pelo contacto de \emph{quantidade}, mas póde movel-o pelo contacto de \emph{força}, emquanto communica ao corpo o ser e a actividade.\par Se a alma humana á \emph{essencial} e \emph{extensivamente} simples, é claro que ella é \emph{indivisivel}. --- Nem se diga que á nossa alma, se não pode ser divisivel \emph{por si}, pode-o ser \emph{por causa do sujeito}, a que está unida. Por quanto uma forma é divisivel \emph{por causa do sujeito}, a que está unida, quando \emph{cada parte}, ainda a mais pequenina, d'esse sujeito tem a \emph{mesma natureza} ou \emph{estructura} do \emph{todo}, de maneira que a forma á por si indifferente tanto para o todo como para uma parte, isto é, pode estar no todo e em cada parte (assim é divisivel a forma do \emph{ferro}); mas não o pode ser, quando o sujeito é um organismo, composto de partes diversas entre si, de maneira que a forma não é e não pode ser indifferente para o todo como para uma parte, e por isso, dividido o sujeito em partes, a forma não se divide e não permanece em cada parte, mas separa-se ou acaba. Ora a nossa alma (como tambem a dos brutos mais perfeitos), por ser forma de um corpo \emph{organico}, não é e não pode ser divisivel por \emph{causa do sujeito}, a. que está unida.}.

%%>>> art-7

\section*{Artigo III.\\ Espiritualidade da alma humana.}
\addcontentsline{toc}{section}{Artigo III. --- Espiritualidade da alma humana.}

%%>>> t32

\tese{32}{Espiritualidade da alma humana.}
\emph{Espiritualidade} é a propriedade, pela qual a alma humana --- não só não é composta de partes, --- mas é \emph{énírinseca} ou \emph{subjectivamente} independente da \pg{39}  materia, tanto nas operações como na existencia. --- Em virtude d' esta propriedade, alma humana diz-se \emph{subsistente} \footnote{Uma substancia póde depender da materia por dois modos: \emph{Intrinseca} ou \emph{subjectivamente}, e \emph{extrinseca} ou \emph{objectivamente}. --- Depende \emph{intrinsecamente}, quando ella, constitue com a materia o principio \emph{adequado} das suas operações. Em tal caso, a cooperação da materia é exigida pela natureza do proprio \emph{sujeito}, e é condição \emph{intrinseca} da actividade d' elle. Esta é a dependencia da alma dos animaes. --- Depende só \emph{extrinsecamente}, quando uma substancia é só por si o principio adequado das suas operações, embora precise do auxilio da outra, que lhe forneça os objectos sobre os quaes possa exercer a sua energia. Em tal caso, a cooperação da materia é só necessaria por causa do \emph{objecto}, e é condição \emph{extrinseca} da actividade do sujeito.\par A alma humana é \emph{intrinsecamente} independente da materia, emquanto não constitue com esta o principio \emph{adequado} das operações \emph{intellectuaes}; porque é só a alma que pensa, e não uma faculdade \emph{organica}, isto é, composta de alma e de corpo. Todavia, a nossa alma depende da materia \emph{ewtrznsecamente}, porque as faculdades sensitivas apresentam os objectos, que depois a alma por si só entende, --- Para a alma ser \emph{espiritual}, basta que seja \emph{intrinsecamente} independente da materia.\par Dissemos que a nossa alma, em virtude da sua \emph{espiritualidade}, é e se diz \emph{subsistente}. Esclaregamos este ponto, que é de grande alcance. -- \emph{Subsistente} é todo o ser que \emph{existe em si e por} si --- que, para existir, não precisa de um outro ser que o sustente, mas que se sustenta a si mesmo; tal é toda e qualquer substancia, que seja \emph{singular} e \emph{individua} (indivisa em si e divisa dos outros seres).--- Um ser \emph{subsistente} pode ser \emph{completo} ou \emph{incompleto}. E' \emph{completo}, quando possue tudo o que é preciso para os seus actos, e não é parte de um composto, mas é perfeito na \emph{propria especie}; tal é um Anjo, um homem, um animal. E \emph{incompleto}, quando não possue por si tudo o que preciso para os seus actos, e por isso é destinado a fazer parte de um composto, e não é perfeito na sua \emph{especie}; tal é a alma humana.\par A alma humana, embora seja \emph{subsistente}, todavia não tem em si e por si tudo o que é necessario para os actos proprios da especie humana, mas precisa do corpo, que lhe seja \emph{comprincipio} nos actos vegetativos e sensitivos, e que tambem lhe subministre a materia, em que a intelligencia possa exercer a sua actividade. Por isso a nossa alma é \emph{incompleta} emquanto á \emph{especie}, \emph{e} é natural e intrinsecamente destinada \& informar o corpo, com que constitue um composto substancial e especificamente completo --- o homem. --- Dizemos que a alma é \emph{incompleta} emquanto á \emph{especie}, e n\&o emquanto á \emph{subsistencia}, porque, na ordem da \emph{subsistencia}, ella é completa, por ter tudo o que é necessario para existir em si e por não carecer de um sujeito que a sustente (Cf. S. Th., \emph{De anima}, art. 1; \emph{De spir}. \emph{creat}. a. 2.\par Do que deixamos dito vê-se a differenca entre a nossa alma e a alma dos brutos, A nossa alma é \emph{incompleta} emquanto á \emph{especie}, mas é \emph{completa} emquanto é \emph{subsistencia}; ao passo quo a alma dos brutos é \emph{incompleta} não só emquanto á \emph{especie}, mas tambem emquanto \& \emph{subsistencia}, porque ella nao existe em si e por si mas existe no composto e pelo composto.}.

%%>>> t33

\tese{33}{Dilferenga entre simplicidade e espiritualidade.}
Ha grande differenga entre \emph{simplicidade} e \emph{espiritualidade}. --- A \emph{simplicidade} \pg{40}  denota negação ou ausencia de partes; a \emph{espiritualidade} significa não só ausencia de partes, mas tambem \emph{Intrinseca} independencia da materia. --- A substancia, que é só \emph{simples}, não é composta de materia, mas está unida á materia e d'ella depende na suas operações e na sua existencia; a substancia \emph{espiritual}, ainda que esteja unida á materia, não depende d'ella \emph{intrinsecamente}, nem nas suas operações, nem na sua existencia. --- A substancia \emph{simples}, por não \emph{subsistir} em si mesma mas no \emph{composto}, acaba com a dissolução do composto; a substancia \emph{espiritual}, por \emph{subsistir} em si mesma, continua a viver fóra da materia, a que está unida \footnote{Alguns escriptores confundiram \& \emph{simplicidade} com a \emph{espiritualidade}, de modo que, segunde elles, o ente, que é simples, é tambem espiritual. D' esta confusão resultaram absurdas consequencias; porque uns, como \emph{Helvetius}, não podendo negar aos animaes um principio \emph{simples} de vida, concederam-lhes o espirito e assim os egualaram ao homem; --- outros, como \emph{Descartes}, não podendo conceder aos animaes o espirito, negaram-lhes tambem um principio \emph{simples} de vida, e fizeram d' elles outras tantas machinas.}.

%%>>> t34

\tese{34}{Methodo a seguir na demonstração da espiritualidade da alma humana.}
A \emph{espiritualidade} não é uma coisa \emph{accidental}, que derive da nossa alma, como o são em geral todas as faculdades; mas é'a propria \emph{essencia} da alma humana. Por isso, assim como a essencia da nossa alma não se póde conhecer \emph{immediatamente} em si mesma, mas só por meio das operações, de que ella é o principio \emph{primeiro}; assim tambem só por meio d'essas operações póde conhecer-se a sua \emph{espiritualidade}.

Este methodo é conforme á razão. Na verdade, entre a essencia de um ser e as suas operações existe a mesma ordem e proporção, que existe entre a causa e o effeito; e por isso da natureza das nossas operações podemos e devemos deduzir a natureza da nossa alma, como da natureza do effeito podemos e devemos deduzir á natureza da causa.

Se, pois, demonstrarmos que em nós ha operações independentes \emph{intrinsecamente} da materia, deveremos concluir que tambem a nossa alma é \emph{intrinsecamente} independente da materia, e por isso \emph{espiritual} \footnote{O methodo que seguimos para demonstrar a \emph{espiritualidade} da alma compõe-se de duas partes: da observação dos \emph{factos} e da applicação dos principios racionaes. Os factos são as nossas operações \emph{intellectuaes}; os principios racionaes resumem-se no principio de \emph{causalidade}: \emph{todo o effeito exige} uma causa \emph{proporcionada}. --- Os \emph{factos} são as nossas operações \emph{intellectuaes}. Na verdade, a alma exerce uma triplice serie de operações : \emph{vegetativas}, \emph{sensitivas} e \emph{intellectuaes}; mas não do mesmo modo. Para exercer as operações \emph{vegetativas}\pgn{41} e \emph{sensitivas}, \emph{a} alma precisa do concurso dos orgãos; ao passo que, para exercer as funcções da vida \emph{intellectual}, a alma não precisa do auxilio dos orgãos. Por isso, se as operações \emph{vegetativas} e \emph{sensitivas} demonstram a \emph{simplicidade} da alma, as \emph{intellectuaes} manifestam a sua \emph{espiritualidade}.\par Quando dizemos que a alma, nas operações \emph{intellectuaes}, opéra sem o auxilio dos orgãos, não queremos dizer que a alma pode pensar regularmente sem o cerebro ou com um cerebro perturbado, mas queremos dizer que não é o orgão material que pensa, mas é só a alma. Por outras palavras: a alma, como dissemos, nas operações \emph{intellectuaes}, \emph{é intrinsecamente} independente dos orgãos, porque é só ella que pensa, embora \emph{extrinsecamente} dependa dos orgãos, que fornecem os materiaes do pensamento.\par As operações \emph{intellectuaes} podem considerar-se \emph{subjectiva} e \emph{objectivamente}. Consideradas \emph{subjectivamente}, são modificações da intelligencia, e não são conhecidas de um modo claro e completo. Consideradas \emph{objectivamente}, isto é, no \emph{objecto} que attingem, são mais accessiveis á investigação; porque pela natureza do objecto conhecemos a natureza das operações e por isso das faculdades e da essencia, pois um objecto \emph{espiritual} só póde ser attingido por uma operação \emph{espiritual}, \emph{e} uma operação \emph{espiritual}, só póde emanar de uma faculdade \emph{espiritual}, e esta só de uma essencia \emph{espiritual}.\par A legitimidade do methodo, que empregamos, é reconhecida pelos mais insuspeitos escriptores. Citaremos as palavras de \emph{B/cAner}, \emph{Vogt} e Wundt.\par \emph{Büchner} diz: « A theoria positivista é obrigada a confessar que o \emph{effeito deve corresponder á causa}, e que por isso effeitos complicados suppõem complicadas combinações da materia » (\emph{Matiére et force}, pag. 218).\par \emph{Vogt}: « E' tambem necessario que a \emph{funcção seja proporcionada d organisação}, \emph{e medida por esta} » (\emph{Lecons sur Ühomme}, pag. 12).\par Wundt: « Não podemos medir \emph{directamente} nem as causas productoras dos phenomenos, nem as forças productoras dos movimentos, \emph{mas podemos medil-as pelos seus effeitos} » (Cit. por Ribot, \emph{Psycol}. \emph{allem}. , pag. 922).\par De resto, este methodo emprega-se constantemente nas sciencias naturaes. Assim a \emph{chimica}, estudando e analysando as propriedades e as operações dos corpos, chega ao conhecimento da natureza intima d' estes.\par Os \emph{positivistas} sustentam que o methodo, que empregamos, não pode produzir certeza e sciencia; porque, dizem, a conclusão --- a \emph{alma é espiritual} --- não póde ser verificada pela experiencia sensivel, e toda a conclusão que não é verificada pela experiencia sensivel é incerta.\par No decurso d' este compendio respondemos mais de uma vez a esta e similhantes objecções. Só accrescentaremos que o argumento dos \emph{positivistas} se\pgn{42}póde retorquir contra elles, pois que podemos dizer: Todo o que-não é verificado pela experiencia não é certo; ora o principio dos positivistas ---\emph{toda} a \emph{conclusão deve ser verificada pela ewcperienciu} --- não foi verificado pela experiencia. Logo este principio não é certo.\par Descartes e os seus sequazes opinam que, para se concluir que a nossa alma é \emph{espiritual}, basta provar que ella é \emph{inextensa} e \emph{indivisivel}. --- Mas esta opinião é falsa. Do facto de a nossa alma ser inextensa e indivisivel só pode concluir-se que ella é \emph{simples}, mas não que é \emph{espiritual}; porque, repetimos, a \emph{simplicidade} denota ausencia de partes e de materia, mas a \emph{espiritualidade} significa ?\emph{ntrinseca} independencia da materia. ---- Os cartesianos partem do falso supposto de que a \emph{espiritualidade} se identifica com a \emph{simplicidade}.\par S. Thomaz trata da \emph{espiritualidade} da alma humana sobretudo na \emph{Summa Theol}, l p., q. 75, a. 2., e nas \emph{Questões Disput}., \emph{De Anima}, a. 1.}.

35. A \textbf{espiritualidade} da \textbf{alma} humana \textbf{demonstra-se pelas operações} da. \textbf{intelligencia}.

á) A nossa intelligencia conhece objectos meramente \emph{immateriaes} ou \emph{espirituaes}: Deus, a \emph{virtude}, \emph{á justiça}, a \emph{ordem}, a \emph{relação} entre a causa e o effeito, ete. Ora, se a nossa intelligencia fosse uma faculdade organica ou dependente da materia, não poderia conhecer similhantes objectos; porque, devendo entre a faculdade e o seu objecto existir uma proporção, uma faculdade organica ou dependente da materia só \pg{42}  póde perceber objectos materiaes, extensos e concretos, mas \textsc{Deus}, a \emph{virtude}, \emph{á justiça}, etc. , são objectos espirituaes, inextensos e independentes da materia. Logo a nossa intelligencia é uma faculdade espiritual. Mas a intelligencia deriva da alma, e um effeito não póde ser superior ao principio de que emana. Logo a alma humana é \emph{espiritual} \footnote{Este argumento, adduzido em favor da \emph{espiritualidade} da alma, é convincente. Uma faculdade organica só póde perceber objectos que, pelas suas propriedades sensiveis, a impressionem. Isto é claro. Ora os objectos, que a intelligencia percebe, não são dotados de propriedades sensiveis, que impressionem a faculdade cognitiva; pois \textsc{Deus}, a \emph{justiça}, a \emph{virtude}, ete., não têm côres para attrahir os nossos olhos, nem harmonia para ferir os nossos ouvidos, nem sabor para interessar a nossa lingua. Logo devemos concluir que --- ou aquelles objectos não são percebidos por nenhuma das nossas faculdades, --- ou que são percebidos por uma faculdade espiritual. Mas não podemos dizer que aquelles objectos não são: percebidos por nenhuma das nossas faculdades; porque isto seria contrario no testemunho da consciencia, Logo são percebidos por uma faculdade \emph{espiritual}.}.

b) A \emph{espiritualidade} da alma demonstra-se não só pelo conhecimento que temos dos objectos \emph{immateriaes} ou \emph{espirituaes}, mas tambem pelo \emph{modo} por que percebemos os objectos \emph{materiaes}. Na verdade, ao passo que os nossos sentidos param na superficie dos corpos e só percebem as suas qualidades sensiveis, a intelligencia, penetrando no intimo dos objectos materiaes, percebe a essencia d'elles, não \emph{material}, \emph{singular} e \emph{contingente}, propria de um individuo, mas \emph{immaterial}, \emph{universal} e \emph{necessaria}, commum a todos os individuos da mesma especie. Ora, se a intelligencia fosse uma faculdade organica e dependente da materia, não poderia conhecer a essencia 2\emph{mmaterial}, \emph{universal} e \emph{necessaria} das coisas corporeas; porque uma faculdade organica só attinge os objectos que nella fazem impressão, e que por isso são \emph{materiaes}. Logo a intelligencia \pg{43}  é uma faculdade espiritual. Se a intelligencia é espiritual, tambem a alma deve ser e é \emph{espiritual} \footnote{Nem este argumento é menos convincente. A nossa intelligencia tem ideas \emph{universues}, formadas pela abstracção. Não só conhecemos este on aquelle homem individuo, mas tambem a humanidade em geral, isto é, a humanidade considerada nos seus elementos essenciaes (abstracta das propriedades particulares, que á tornam concreta e singular), e por isso commum a todos os homens. Sem ideas universaes, não é possivel a sciencia, nem o juizo, nem o raciocinio. A intelligencia, pois, tem ideas universaes. Ora, se a intelligencia fosse uma faculdade organica, não poderia apprehender ideas \emph{universaes}. Com effeito, uma faculdade organica só apprehende o que lhe é accessivel o que a póde impressionar, e por isso o que existe na ordem physica e real Mas o universal, como \emph{tal}, isto é, abstracto das condições individuaes, não existe na ordem physica e real; porque no mundo não existe a essencia, separada das condições individuaes e concretas, mas existem individuos dotados d'aquella essencia. Logo a intelligencia é uma faculdade inorganica e espiritual; visto que só uma faculdade inteiramente independente da materia pode attingir objectos immateriaes, pois a faculdade deve ser proporcionada ao objecto. Logo a nossa alma é \emph{espiritual}.}.

o) A intelligencia humana não só percebe os objectos \emph{immateriaes} e os \emph{materiaes} de um modo \emph{immaterial}, mas tambem reflecte sobre si mesma, vé o seu acto, e affirma a propria identidade e permanencia atravez da multidão das ideas, que se suecedem constantemente; de modo que ella é, ao mesmo tempo, faculdade cognitiva e objecto conhecido. Mas, se a intelligencia fosse uma faculdade organica, não poderia reflectir sobre si mesma; porque uma faculdade organica é composta de partes extensas, e uma parte extensa póde reflectir sobre uma outra parte, mas nunca sobre si mesma. Logo a intelligencia humana é uma faculdade espiritual. Mas uma faculdade espiritual só pode derivar de uma essencia espiritual.

\emph{Z} Logo a alma humana é \emph{espiritual} \footnote{Tambem este argumento prova convincentemente a proposição. --- A nossa alma reflecte sobre o proprio pensamento, pensa que pensa, conhece-se a si mesma como \emph{sujeito} das suas operações, e percebe a sua essencia, distinguindo-a do seu acto. D'este modo, a alma é, ao mesmo tempo, sujeito e objecto do seu pensamento, e a sua operação é eminentemente \emph{immanente}. Ora, se a alma pensasse por meio de um orgão, não poderia reflectir sobre si mesma, nem conhecer-se como \emph{sujeito} das proprias operações. --- Não poderia reflectir sobre si mesma, porque seria necessario que o orgão se objeotivasse, isto é, se tornasse objecto de si mesmo. Ma isto é impossivel; porquanto no corpo uma parte póde obrar sobre outra parte, mas nunca reflectir Sobre si mesma. --- Nem poderia conhecer-se a si mesma como \emph{sujeito} das suas operações; porque, para isso, deveria perceber a relação de \emph{causalidade}; mas esta relação, por ser uma coisa \emph{abstracta}, não póde ser percebida por\pgn{44}uma faculdade organica. Logo a nossa alma não pensa por meio de um orgão material mas é \emph{intrinsecamente} independente da materia, é \emph{espiritual},}.

%%>>> t36

\pg{41} \pg{44} \tese{36}{A espiritualidade da alma humana demonstra-se pelas operações da vontade.}

a) A nossa vontade deseja bens immateriaes e espirituaes, como são a \emph{virtude}, a \emph{sciencia}, a \emph{posse} de \textsc{Deus}, ete., e nestes bens encontra a sua perfeição e a sua felicidade. Se deseja bens espirituaes, a nossa vontade é espiritual; não só porque todo o ser só tende para o bem que-Ihe é similhante, mas tambem porque bens espirituaes ou immateriaes não podem impressionar uma faculdade organica ou dependente da materia. Se a vontade é espiritual, tambem a alma é \emph{espiritual}; porque a vontade, como a intelligencia, deriva da essencia da alma e nunca um effeito póde exceder em perfe: Co á causa, de que deriva \footnote{As faculdades \emph{appetitivas} de um ser estão em proporção com as suas faculdades \emph{perceptivas}; porque um ser tende para o objecto que conhece. Se as faculdades \emph{perceptivas} attingem objectos materiaes, as \emph{appetitivas} tendem para objectos materiaes; como tambem, se as faculdades \emph{perceptivas} attingem objectos \emph{espirituaes}, as \emph{appetitivas} tendem para objectos \emph{espirituaes}. A nossa intelligenecia apprehende objectos espirituaes; e por isso a nossa vontade tende para objectos espirituaes. Se a intelligencia, pelo facto de apprehender objectos \emph{espirituaes}, \emph{é espiritual}; tambem a vontade, pelo facto de tender para objectos \emph{espirituaes}, \emph{é espiritual}, porque, como dissemos, toda a faculdade tende para os objectos que lhe são similhantes.}.

b) Se a nossa vontade fosse uma faculdade organica e dependente da materia nas suas operações, não poderia levantar-se sobre a materia e dominal-a; pois repugna que o instrumento domine o agente pelo qual é movido. Ora a nossa vontade effectivamente rege e domina a materia, sujeitando- á ás dores, aos tormentos e á morte, apezar de toda a resistencia que oppõe o appetite sensitivo. Logo a nossa vontade é independente da materia, e por isso é \emph{espiritual}. Logo a alma bumana é \emph{espiritual} \footnote{Note-se a differenca entre as operações do appetite do animal e as da vontade do homem. --- O appetite do animal segue \emph{necessariamente} as impressões organicas; a vontade do homem póde resistir e muitas vezes resiste ás impressões organicas. --- O appetite instinctivo do animal póde ser subjugado por uma força externa; a vontade do homem póde resistir e muitas vezes resiste a todas as pressões dos agentes externos, de modo que, se a materia cede, a vontade triumpha. Ora, a vontade do homem não poderia contrariar e dominar a materia, se não fosse inteiramente independente da propria materia, isto é, se não fosse espiritual.}.

\emph{Z}

%%>>> art-8

\section*{Artigo IV.}
\addcontentsline{toc}{section}{Artigo IV.}

{\small\noindent\textsc{Summario}: --- União da alma com o corpo e suas consequencias. --- Unidade da alma humana. --- Séde da alma humana.\par}

\pg{45} A \textbf{espiritualidade da alma humana}

e o \textbf{materialismo}.

%%>>> t37

\tese{37}{Doutrina do materialismo.}
O \emph{materialismo} pretende que o pensamento não é o producto de uma faculdade inorganica, mas é apenas uma secreção, um movimento, uma funeção do cerebro, porque a materia por si é capaz de pensar.

Provaremos que a materia é absolutamente incapaz de pensar, e que por isso o pensamento não é uma secreção, nem um movimento, nem uma funceção do cerebro \footnote{As asserções do \emph{materialismo} foram já indirectamente refutadas no artigo precedente, em que demonstrámos que o pensamento é o producto de uma faculdade espiritual ou independente da materia. Mas, tratando-se de um systema, que, em nome da sciencia, espalha as mais absurdas e pernieiosas doutrinas, julgamos necessario refutal-o directamente, ainda que haja, por vezes, alguma repetição de certos principios fundamentaes.}.

%%>>> t38

\tese{38}{Repugna que o pensamento seja um producto da materia.}
Repugna que um effeito seja dotado de attributos diametralmente oppostos, ou superiores, aos attributos de que é dotada a sua causa; porque esse effeito não teria uma causa proporcionada. Ora o pensamento é dotado de attributos oppostos, e superiores, aos que possue a materia. Na verdade, o pensamento, como temos provado, é abstracto e incorporeo, não tem figura nem peso, e prescinde das cireumstancias ou condições de tempo e de logar; ao passo que a materia é necessariamente concreta e sensivel, tem figura e peso, e está determinada pelas eireumstaneias de tempo e de logar. Logo repugna que o pensamento seja um producto da materia \footnote{\emph{Locke}, embora admittisse uma distincção entre a alma ea materia, todavia opinou que só pela razão não podemos demonstrar se o principio pensante no homem a uma substancia distincta do corpo, ou se é a propria materia que é dotada da força de pensar. E' a celebre \emph{duvida} de \emph{Locke}.\par No texto demonstramos que a materia, em quanto fôr materia, é incapaz de pensar. Materia e pensamento são duas coisas incompativeis.}.

%%>>> t39

\tese{39}{O pensamento não a uma secreção do cerebro.}
Tudo o que um orgão segrega, é material, divisivel, ponderavel. Ora o pensamento --- é \emph{Immaterial}, porque apprehende o universal, o espiritual; --- é \emph{Indivisivel}, porque repugna que uma idea abstracta, como a da \emph{virtude}, da \emph{justiça}, possa dividir-se em metades, em quartos, etc.; --- é \emph{zmponderavel}, porque é absurdo \pg{46}  dizer que uma idea pesa, p. ex. , vinte ou quarenta grammas. Logo o pensamento não é uma \emph{secreção} do cerebro \footnote{\emph{Vogt} tornou-se celebre por esta phrase: « O pensamento está para com o cerebro na mesma relação em que a bilis está para com o figado. n\par A opinião de \emph{Vogt}, que é apenas uma repetição da de \emph{Cabanis} (\emph{Rapports du physique} ecc.) é absurda. Na verdade, a producção de uma coisa está em harmonia coma sua natureza; de modo que duas coisas serão diversamente produzidas, se diversas forem as suas naturezas. Ora a natureza da bilis é essencialmente diversa da do pensamento. Por quanto, --- a bilis é uma substancia visivel, \emph{o} pensamento é uma coisa invisivel, --- a bilis é uma secreção, porque se expelle do figado, mas o pensamento não é uma secreção, porque não se expelle da faculdade cognitiva, mas fica nella.\par De resto, os mais insuspeitos escriptores, como \emph{Büchner}, \emph{Maudsley}, \emph{Lewes}, \emph{Lange}, declararam que aquella comparação foi mal escolhida.\par Citaremos apenas as palavras de \emph{B/üchner}: « O mais minucioso exame não permitte descobrir uma analogia real entre a secreção da bilis e o processo pelo qual se forma no cerebro o pensamento. A bilis é uma substancia palpavel, ponderavel e visivel, e, além d'isso, é materia excrementicia, que se expelle do corpo; pelo contrario, o pensamento, longe de ser materia excrementicia, constitue uma actividade ou um movimento das substancias ou das combinações das substancias dispostas no cerebro por um modo determinado » (\emph{Force et matiére}, pag. 892).}.

%%>>> t40

\tese{40}{O pensamento não: é um movimento do cerebro.}

a) O movimento de um ente material é um phenomeno exterior ou transeunte, é a modificação de uma coisa extensa, figurada, situada no espaço e determinada pelo tempo; ao passo que o pensamento é uma coisa interior ou immanente, não tem extensão nem figura, prescinde do tempo e do espaço.

b) De mais, se o pensamento fosse um movimento material, a mesma faculdade, que percebe o pensamento, deveria perceber o movimento; e vice-versa. Mas não é assim. Pela consciencia, percebemos o pensamento, mas não o movimento; e, pelos sentidos, percebemos o movimento, mas não o pensamento.

c) Finalmente, se o pensamento fosse um movimento da materia, o pensamento deveria estar sujeito a todas as leis, a que esti sujeito o movimento. Todavia não é assim. A substancia, que está em movimento, deve tocar todos os pontos do espaço, que percorre; mas o principio pensante, quando pensa, passa de repente de um ponto para outro mais remoto, sem comtudo tocar os pontos intermedios.

Logo o pensamento não é um movimento do cerebro \footnote{Se o pensamento fosse apenas um movimento das cellulas cerebraes, o juizo seria impossivel. Porquanto, quando julgamos, unimos ou separamos duas ideas, --- o sujeito e o predicado. Mas estas duas ideas não se podem unir ou separar sem uma previa comparação, e esta comparação deve\pgn{47}ser feita pela mesma faculdade, que perceba uma e outra idea. Assim, quando dizemos: \emph{e ente não a o nada}, se as duas ideas « \emph{ente} e \emph{nada} » fossem dois movimentos, seria impossivel o juizo. Na verdade, estes dois movimentos teriam logar na mesma cellula ou em diversas cellulas. Se tivessem logar na mesma cellula, o juizo seria impossivel por falta de termos; pois o \emph{ser} e o \emph{nada}, sendo oppostos entre si, destruir-seiam, porque dois movimentos conbrarios no mesmo sujeito destroem-se. Se as duas ideas residissem em duas cellulas, tambem neste caso o juizo seria impossivel, porque faltaria a comparação, pois não é a mesma faculdade que percebe o predicado e o sujeito\par \emph{Paulo Janet} diz: « Um movimento póde ser rectilineo, cireular, ou em espiral: o que é um pensamento em espiral, circular ou rectilineo?!.. O meu pensamento é claro ou obscuro, verdadeiro ou falso: o que é um movimento claro ou obscuro, verdadeiro ou falso?! Numa palavra, um movimento pensante envolve contradiecao » (\emph{Traité élément}, \emph{de Philosophie}, pag. 826).\par \emph{J}. \emph{Rosenthal} escreve: « Podemos comprehender como um movimento produza outro movimento, mas desconhecemos inteiramente como um movimento possa produzir uma percepção » (\emph{Les muscles et les nerfs}).\par Foram \emph{AMolleschott}, \emph{Büchner}, \emph{Herzen}, \emph{Richet}, e outros materialistas, que reduziram o pensamento a um movimento das moleculas cerebraes.\par Os \emph{materialistas} objectam: Quando pensamos, as moleculas cerebraes movem-se. Logo o pensamento é um movimento do cerebro.\par Concedemos o facto, mas negamos a consequencia, que d'elle tiram os materialistas. Do facto, apontado pelos adversarios, só pode deduzir-se que o pensamento anda acompanhado pelo movimento das moleculas cerebraes, mas não se deduz que o pensamento é o proprio movimento. Como provaremos d' aqui a pouco, a actividade da alma \emph{agora} não se exerce sem o concurso dos sentidos internos; e, como os sentidos internos têm o seu orgão em certas regiões do cerebro, segue-se que, exercendo a alma as operações intellectuaes, tambem o cerebro deve entrar em movimento.\par O proprio \emph{Taine} confessa que o pensamento e o movimento não podem identificar-se um com outro, porque irreductiveis (\emph{De l'intellig}. t. I, pag. 885).}.

%%>>> t41

\pg{47} \tese{41}{O pensamento não é uma funcção do cerebro.}

a) Toda a funcção do cerebro é sensivel, concreta e determinada pelo tempo e pelo espaço. Ora o pensamento, como dissemos, escapa á percepção dos sentidos, é universal e pre-

4 seinde de todo o tempo e de todo o espaço.

b) Se o pensamento fosse uma funeção do cerebro, entre a faculdade pensante e o proprio cerebro deveria existir \emph{sem pre} uma equação completa e rigorosa. Ora a experiencia demonstra que similbante equação não existe \emph{sempre}.

Logo o pensamento não é uma funeção do cerebro \footnote{Os \emph{materialistas}, para provar que o pensamento é uma funeção do cerebro, recorrem á \emph{anatomia} e á \emph{physiologia}. Na opinião d'elles, a \emph{anatomia} diz que entre o pensamento e o cerebro existe \emph{sempre} uma equação \emph{completa} e \emph{rigorosa}, e á \emph{physiologia} diz que o pensamento é uma \emph{funcção} do cerebro.\par Respondemos que, embora entre o cerebro e o pensamento existisse \emph{sempre} uma rigorosa e perfeita equação, isto só provaria que o cerebro é uma\pgn{48}\emph{condição} do pensamento, e não que é \emph{causa}. Mas é falso que á \emph{anatomia} ensine que entre o cerebro e o pensamento existe aquella equação; como é falso que a \emph{physiologia} ensine que o pensamento é uma funceção do cerebro.\par Em defeza d'estas nossas proposições, provaremos primeiramente que a \emph{anatomia} não demonstra existir entre o cerebro e o pensamento uma \emph{rigorosa} equação, e que a \emph{p/jsiologia} não ensina ser o cerebro o orgão do pensamento; em seguida, provaremos que o cerebro é apenas o orgão das \emph{condições} do pensamento, e finalmente faremos a critica d'algumas asserções dos mais famigerados materialistas.\par O estudo \emph{anato:zico} do cerebro não é recente, pois vem de edades remotas; mas o estudo das suas relações com o pensamento é moderno e principiou com a reacção, que tinha por fim recompor o homem e restituir-lhe a unidade pessoal, que \emph{Descartes} tão inconsideradamente tinha quebrado.\par a) \emph{Gall} foi quem primeiro tentou resolver o problema das relações entre o cerebro e ó pensamento. Defensor acerrimo da unidade pessoal do homem, o naturalista allemão indicou o caminho, que um sabio deve seguir nas suas investigações ácerca do homem, mas não consoguiu o seu fim. Defendeu a espiritualidade da alma humana, mas inventou um systema, que devia subministrar muitas armas aos materialistas.\par O \emph{phrenologismo} foi recebido com enthusiasmo; pois uma simples observação do craneo permittia a todos o conhecer as faculdades intellectuaes predominantes no individuo. Segundo \emph{Gal}/, no cerebro distinguem-se muitas e diversas proeminenoias. Cada uma d'ellas é orgão de uma faculdade primitiva e superior. E, como a forma do craneo corresponde perfeitamente á forma do cerebro, legitimamente nas saliencias do craneo podemos reconhecer e distinguir as faculdades da alma. --- .\emph{Spurzheim}, discipulo de \emph{Gall}, encontrou mais proeminencias no cerebro, e deu ás faculdades da alma uma disposição ou localisação diversa da do seu mestre.\par Esta theoria não tinha uma base solida. \emph{Gall} e \emph{Spursheim} dividiram o cerebro e localisaram as faculdades a capricho e sem criterio; por isso o \emph{phrenologismo} encontrou opposição entre os sabios, e foi condemnado pela razão, pela sciencia e pela experiencia. --- Foi condemnado pela razão; porque faculdades espirituaes, como são a intelligencia e a vontade, não podem residir em orgãos materiaes, --- Foi condemnado pela sciencia; porque a anatomia demonstra que a conformação exterior do craneo não corresponde á sua conformação interior e por isso á forma do cerebro. --- Foi condemnado pela experiencia; porque no craneo de \emph{Laplace} foi encontrado o orgão da estupidez, no craneo do famoso assassino \emph{F/esch} foi encontrado o orgão da benevolencia e no do celebro ladrão \emph{Lacenaire} o orgão da liberalidade.\par b) Os materialistas, não encontrando no craneo dos vivos razões para a defeza das suas doutrinas, foram proeural-as no craneo dos mortos. Mediram a capacidade do craneo dos mortos; porque d'este modo, diziam, temos a quantidade exacta da massa cerebral, e por isso a quantidade da intelligencia; pois grande volume do encephalo denota necessariamente grande\pgn{49}talento. --- \emph{Büchner} recorre á auctoridade dos chapelleiros, e diz que a quotidiana experiencia d'estes officiaes ensina que a classe culta precisa de chapeus maiores que a classe ignorante.\par Se fosse verdade que uma capacidade maior do craneo denota sempre maior intelligencia, os craneos dos antepassados deviam ter uma capacidade inferior a dos nossos; pois os materialistas dizem que os antepassados eram dotados de uma intelligencia mediocre, muito inferior a dos homens da nossa epoca. Todavia a experiencia mostra o contrario. Encontraram-se \emph{seis} craneos, nas grutas de Lozére, na Franca, pertencentes a uma epoca remotissima, a que deram o nome de \emph{epoca da pedra polida}. Pois bem, cada um d'esses craneos tinha uma capacidade media de 1606 centimetros cubicos; ao passo que da experiencia, feita em 67 craneos de individuos parisienses, do ultimo seculo, resultou que estes craneos tinham uma capacidade media de 15596. c., isto é, 4c. c. menos que os craneos da \emph{pedra polida}.\par Mas façamos uma comparação entre os craneos dos individuos de diversas nações na epoca presente. A capacidade do craneo dos Esquimós é de 1546c.c., a dos Neo-Zelandezes é de 1497c.c., a dos Japonezes de 1486.o., a dos Europeus de 1467e.c. Se fosse verdadeira a opinião dos materialistas, os Esquimós seriam os mais intelligentes do mundo, e os Europeus só occupariam o quarto grau na escala da intelligencia; o que é falso.\par c) \emph{Topinard} e os seus discipulos ficaram deveras embaracados á vista d'este resultado, como de um phenomeno imprevisto; mas não se \textbf{deram} por vencidos. Mudaram apenas o campo da contenda, e disseram que não é a capacidade do craneo que se deve medir, mas é o proprio cerebro que se deve pesar; pois o peso do cerebro augmenta com o grau do desenvolvimento intellectual; quanto mais desenvolvido fór o cerebro, tanto mais poderosa será a intelligencia.\par Este recurso não salva o materialismo. Appellemos para a \emph{anatomia}, que nos dá os seguintes resultados:\par O cerebro de um elephante pesa em media . . 93000 gr.\par (i doSdoLpHims Goss un ce cer ess Un sr ue D\textasciicircum{} ue rope e 1800! »\par O da baleias cu vasos sa s ter per e a 00058\par Odo homem. .- IA v4 95 wo. 4900) 4\par Se, pois, o desenvolvimento do cerebro fosse proporcionado ao desenvolvimento intellectual, o elephante, o delphim e a baleia seriam mais intelligentes do que o homem. ]\par d) Os materialistas dizem que não se deve attender ao peso \emph{absoluto} do cerebro, mas sim ao peso \emph{relativo}.\par Respondemos que este recurso é perfeitamente inutil e vão. Prescindindo dos outros animaes, que formam a escala zoologica, e considerando só os mammiferos, temos os seguintes resultados. No cavallo, o peso do cerebro está para o peso do corpo como 1 para 400; no macaco, como 1 para 100; no homem, como 1 para 47; no saguim, como 1 para 28; no canario, como 1 para 14. Se, pois, valesse o argumento dos materialistas, o. saguim\par 4 Vor. II.\pgn{50}e o canario deveriam ser mais intelligentes do que o homem, pois o peso do cerebro d'aquelles animaes é \emph{relativamente} maior que o do homem.\par e) Oa materialistas, porém, querem limitar-se ao peso \emph{relativo} do cerebro das raças humanas, e dizem que, na mesma especie humana, o peso do cerebro das diversas raças está em rigorosa proporção com os differentes graus de civilisacüo e por isso de intelligencia.\par Similhante asserção, que, ha pouco tempo, era tida como um axioma scientifico, está hoje abandonada, como contraria á experiencia.\par Eis o peso do cerebro de alguns povos.\par O cerebro dos Escocezes pesa em media. . . . 1417 gr.\par O dos Insulares das Carolinas. . . . . . . . 1402 m\par O dos Esquimós. . . . 4 or ot on n 1398 «\par O dos Inglezes. . a « ROS so eee des 1988 »\par O dos Bavaros! A Lui ee 10. e deo se. ts STS car\par O dos Franceges. 7. "Ug le eba pe e, ABS On\par O'dos Annümitás 19 eo RO eoe e ABES\par O dos Italianos. 10 59 woe eu ae e o 808 nm\par Na opinião, pois, dos materialistas, os Insulares das Carolinas, os Esquimós e os Annamitas deveriam ser mais intelligentes do que os Inglezes, os Bavaros, os Francezes e os Italianos.\par 1) A falsidade do systema materialista torna-se ainda mais evidente, se considerarmos o cerebro no seu desenvolvimento. --- Eis os resultados anatomicos, segundo as experiencias de \emph{Boyd}.\par No nascimento, o cerebro humano pesa. . . . 981 gr.\par Na edade de 8 Mmeges . x suis 0 495»\par Da 3a 6.106268 + 5 « s xm 9 ne Lom sed s 608 m\par De é mezesa 1 anno.. . . s á o: \emph{Tt}»\par De. lia FANnDO8 o x * 7 € [95-10 9-4 + 949 «\par Domos 4e nnnps o. 53 etc sara ea E Mee MES 1097 -\par Do-4 ar (-RBDOS-- s ce] \emph{em eT ed} 78 Wise 1140 »\par DS Tia 14 RnnoB e rs NORD S NES E S SE 1902\par Do la 20 \&nnos s € v9 3 e v. \$3 T t 1894 »\par Do 20 m BO SHHOB eus cum o room DR DEUS 13857 »\par Do 80 a 40 annos e 9 5s s € \$9. e o5 1366 »\par «Dacd0g2DOCRnBboS- é coe v e red ue Ee S 1852.»\par De:sB0lassGÜ mnA08: x1 se s x e mer re T 1948\par Doe 60 a TOMBOS: 115 com re ce e 13815 »\par De TO a 80 ABDOS. 2 ara Of OE T Te Ue C 1289 »n\par Da 80 açdO Annos. . á e 21 pte NE e eh 1984 \textasciicircum{}\par Conhecidos estes resultados, tiramos as seguintes conclusões.\par 1.) Se o cerebro fosse o orgão do pensamento, a criança, que tem um \textbf{anno} de edade, já deveria pensar, porque tem 717 grammas de cerebro, isto é, um volume superior ao volume do macaco já adulto, que, segundo as doutrinas dos materialistas, é dotado de intelligencia.\par 2.) De mais, se o grau da intelligencia estivesse na razão directa do peso\pgn{51}do cerebro, todas as crianças de 12 annos, cujo cerebro pesa 1260 grammas, deveriam ser mais intelligentes do que o celebre \emph{Gambetta}, cujo cerebro apenas pesava 1160 grammas, isto é, 140 gr. menos que a media.\par 8.) Finalmente, se do peso do cerebro dependesse o grau da intelligencia, então o homem, depois dos 40 annos, quando o peso do cerebro principia a declinar, deveria declinar nas faculdades intellectuaes. Mas não é assim; porque todos sabem que, geralmente, depois dos 40 annos principia a edade da reflexão, do juizo e do conselho.\par Verdade é esta tão evidente, que um dos mestres materialistas, \emph{Broca}, apertado pelas objecções de \emph{Gratiolet}, chegou a fazer a seguinte declaração, que está em contradicção aberta com as doutrinas materialistas: » Ninguem pretendeu, nem pode pretender, que exista uma relação \emph{absoluta} entre o desenvolvimento da intelligencia e o volume do cerebro. Protestei sempre contra este absurdo. O sabio nem sequer póde pensar em que, medindo o cerebro, meça a intelligencia n.\par Todos estes desmentidos da \emph{anatomia} deviam ser mais que sufficientes para obstar a que os materialistas apregoassem as suas doutrinas em nome da sciencia. Mas não succedeu assim. Obrigados a reconhecer que nem nas proeminencias e na cavidade do craneo, nem no peso \emph{absoluto} ou \emph{relativo} do cerebro encontravam uma base para o seu edificio, os materialistas disseram que, em fim, o volume do cerebro é um dos factores da intelligencia, mas não é \emph{unico}; pois o grau mais ou menos elevado da intelligencia depende da symetria dos hemispherios cerebraes, da quantidade da substancia cinzenta e de phosphoro, e da grandeza dos lóbos frontaes. .\par Examinemos estas affirmações.\par g) Foi o naturalista \emph{Bichat}, quem fez depender o maior ou menor desenyolvimento da intelligencia da symetria, mais ou menos perfeita, dos hemispherios cerebraes.\par Mas esta opinião morreu com o proprio auctor. Na autopsia feita ao cadaver de \emph{Bichat}, um dos hemispherios cerebraes foi encontrado atrophiado; e todavia esse anatomista era intelligentissimo. Em seguida, a sciencia demonstrou que o atrophiamento parcial ou total de um só hemispherio póde deixar e muitas vezes deixa intacta a intelligencia.\par h) Appellando para a substancia \emph{cinzenta}, os materialistas disseram que os que têm mais substancia \emph{cinzenta}, tem mais intelligencia. E, como a substancia \emph{cinsenta} forma a camada cortical do cerebro, sustentaram que o grau da intelligencia será mais ou menos elevado, segundo forem mais ou menos salientes as \emph{circumvoluções}.\par Não gastaremos muitas palavras para refutar esta doutrina. Sô diremos que, se devessemos acreditar nas palavras dos materialistas, então mais intelligentes do que todos os outros animaes, mais intelligentes do que o proprio homem seriam o carneiro e o... jumento; porque as \emph{circumvoluções} do cerebro d' estes dois animaes são salientissimas, e, pelo contrario, o mais estupido de todos os animaes seria o castor, que tem o cerebro liso.\par i) Outros materialistas, como \emph{Moleschott}, \emph{Couerbe}, \emph{Feuerbach}, attribuem o maior ou menor desenvolvimento da intelligencia á maior ou menor quantidade\pgn{52} de \emph{pAhosphoro} contido no cerebro; e dizem que a abundancia de \emph{phosphoro} no cerebro produz gehio; o seu excesso, loucura; e a sua ausencia ou escassez, idiotismo. :3\par Esta opinião foi refutada por muitos sabios, e especialmente por \emph{Fréêmy}. As analyses chimicas, realizadas no cerebro de muitos loucos, mostraram que elle não tem mais \emph{p/Aiosphoro} do que o dos outros homens. --- De mais, o cerebro dos peixes e o do animal suino tém mais phosphoro, que o do homem; e, todavia, ninguem dirá que esses animaes são mais intelligentes do que o homem. --- Finalmente, o phosphoro desenvolve-se no cerebro dos animaes depois da sua morte, e por isso não póde ser a causa da intelligencia.\par j) Outros materialistas dizem que a intelligencia está localisada nos \emph{lóbos frontaes}. --- \emph{Topinard} diz: « Nos lóbos frontaes residem as mais altas faculdades. Lóbos volumosos, e por isso fronte espaçosa, revelam grande intelligencia; lóbos atrophiados, e por isso fronte baixa, são indicio de uma intelligencia mediocre ou nulla ».\par Respondemos que o desenvolvimento dos lóbos frontaes não é indicio seguro do desenvolvimento da intelligencia. Se em alguns idiotas esses lóbos estavam atrophiados, em outros conservavam o seu estado normal. Se Walter \emph{Scott} e \emph{Cuvier} possuiam uma fronte alta, o celebre \emph{Lacépéde} e o General \emph{Lamarque} eram dotados de uma fronte baixa.\par k) Outros materialistas, como \emph{ANewmann de Leipzig}, collocaram a intelligencia nos lóbos occipitaes, porque no cerebro de muitos loucos esses lóbos estavam mais atrophiados do que os frontaes.\par Esta opinião, como a antecedente, está desmentida pela moderna theoria das \emph{localisagóes}, --- theoria que tem por auctores dois physiologistas allemães, \emph{Fritsch} e \emph{Hitzig}. Estes sabios demonstraram até á evidencia que os \emph{centros motores} residem quasi exclusivamente no lóbo parietal e na metade posterior do lóbo frontal, e que a sensibilidade reside na parte posterior dos hemispherios e especialmente nos lóbos occipital e temporo-esphenoidal. D'este modo todo o cerebro só é orgio do movimento e da sensibilidade. A intelligencia, pois, não tem orgão, e por isso não é uma funcção do cerebro. --- E' verdade que os materialistas dão aos \emph{centros motores} o nome de \emph{centros psychico-motores}, e fazem d'esses centros o orgão da intelligencia; mas esta pretencüo foi eliminada pela sciencia, que nas camadas corticaes reconheceu o unico e exclusivo officio de \emph{centros motores}.\par \emph{A anatomia}, pois, longe de favorecer o materialismo, refuta-o.\par Coneluiremos com Waitz; « Duvidamos que haja um só individuo, dotado de alguns conhecimentos physiologicos, que se persuada de que o espirito se avalie pelo peso do cerebro, se meça pela grandeza do craneo, pelo desenvolvimento comparativo das differentes partes da massa encephalica, e de que as innumeraveis aptidões ou disposições intellectuaes, moraes e sentimentaes se desenhem nas protuberancias da superficie do craneo ».\par Os materialistas recorrem tambem á \emph{pysiologia}, e, em nome d'estasciencia, proelamam que o cerebro é orgão do pensamento, pois que está averiguado que o cerebro \emph{trabalha}, quando o homem pensa.\par Vejamos se verdadeiramente a \emph{physiologia} dá razão aos materialistas.\par E' certo que o cerebro \emph{trabalha}, quando o homem pensa.\par Agora pergunta-se: O cerebro será \emph{o orgão}, \& \emph{causa efficiente} do \textbf{pensamento}, ou apenas uma \emph{condição} ?\par Os materialistas dizem que o cerebro é o \emph{orgão}, a \emph{causa efficiente} do pensamento; nós, pelo contrario, sustentamos que é apenas uma \emph{condição}. ;\par Será verdade que a \emph{plusiologia} ensina ser o cerebro o \emph{orgão} do pensamento? --- Respondemos que não. Temos em nosso favor o testemunho dos celebres \emph{Gavarret} e \emph{Griesinger}, e dos insuspeitos \emph{Du Bois}- \emph{Reymond} e \emph{Tyndall}.\par \emph{a}) \emph{Gavarret} assim se exprime: « Da parte do cerebro ha um augmento de actividade e de combustão, e uma producção de calor; e, ao mesmo tempo, dá-se uma manifestação psychica. Entre aquelle trabalho interior e essa manifestação existe uma coincidencia perfeita; o primeiro é evidentemente uma \emph{condição} da segunda. Mas haverá, além d' esta, uma outra relação entre os dois phenomenos? Existirá uma medida commum para a quantidade de calor que se exhala e um pensamento que se concebe? Emquanto se não encontrar essa medida commum, não temos o direito de affirmar --- \emph{que o trabalho cerebral e a concomitante manifestação psychica só differem na forma}, --- \emph{que estes dois productos têm a mesma natureza}, --- \emph{que o primeiro} é causa \emph{sufficiente do segundo} « (\emph{Les phénoienes physiques de la vie}, 1869).\par \emph{b}) \emph{Griesinger} diz; « Como é que um phenomeno material e physico, que se realiza nas fibras nervosas ou nas cellulas ganglionarias, póde tornar-se uma idea, ou um acto de consciencia? Isto é absolutamente incomprehensivel » (\emph{Traité des maladies mentales}, trad. Doumic. 1865).\par \emph{e}) \emph{Du Bois-Reymond} escreve: « Ainda que conhecessemos intrinsecamente o cerebro, os phenomenos intellectuaes seriam sempre incomprehensiveis.... O mais profundo conhecimento do encephalo só nos manifesta a materia em movimento. Mas nenhuma disposição, nenhum movimento de partes materiaes póde servir-nos de ponte de passagem para o, dominio da\par SEIEN \emph{tem} \textbf{Antelligenoia}. . . Á primeira vista e por uma observação superficial parece que, pelo conhecimento dos factos materiaes no cerebro, poderiamos comprehender certos factos de ordem espiritual...; mas a mais breve reflexão mostra que isto é uma illusão. Apenas podemos obter noções sobre as cireumstancias da vida espiritual, que se tornam manifestas pelas impressões dos sentidos, mas não sobre a natureza da vida espiritual por meio d'essas circumstancias... De modo nenhum se póde comprehender como a consciencia possa nascer da combinação de certo numero de atomos de carbone, hydrogenio, azote, oxygenio, etc. Não é preciso demonstrar a impossibilidade actual e permanente de comprehender os elevados phenomenos espirituaes pela supposta mechaniea dos atomos cerebraes... Os nossos conhecimentos naturaes estão encerrados entre dois limites: de um lado a incapacidade para comprehender a materia e a força, do outro a impossibilidade de deduzirdas condições materiaes os phenomenos espirituaes: estes limites estão determinados irrevogavelmente » (\emph{Sur les bornes de la philos}. \emph{naturelle}).\par \emph{d}) \emph{Tyndall} diz: « E' impossivel conceber a passagem da physica do cerebro ao facto correspondente da consciencia intima das sensações, dos pensamentos, das emoções. Se nos concedessem que um pensamento determinado e uma determinada acção do cerebro são factos simultaneos, não só não teriamos demonstrado a existencia do orgão intellectual, mas nem sequer a existencia de um rudimento visivel d'este orgão. Ainda quando a nossa intelligencia e os nossos sentidos fossem bastante aperfeiçoados, vigorosos e illuminados para deixarem ver e sentir as ultimas particulas do cerebro, e seguil-as nos seus movimentos; ainda quando tivessemos a consciencia dos estados correspondentes do pensamento e do sentimento, estariamos sempre muito longe da solução d'este problema. Qual é a relação entre as operações physicas e os factos da consciencia? O abysmo entre as duas classes de phenomenos é e será sempre insuperavel » (\emph{Les mondes}).\par Do mesmo modo se exprimem \emph{Jaccoud} (\emph{Traité de pathologie interne}), Potain (Art. \emph{Cerveau}, no \emph{Dict}. \emph{encycl}. \emph{des sciences med}.) Ferrier (\emph{Fonctions du cerveaw}) Vierordt (\emph{Le physiologie de l'homme}).\par Estas declarações tiram ao materialismo toda a apparencia scientifica. A physiologia constata o facto, diz que existe uma relação entre o trabalho do cerebro e o pensamento, mas não póde affirmar se essa relação é \emph{causal} ou apenas \emph{condicional}. --- Para a physiologia, o pensamento é incomprehensivel- Na verdade, esta sciencia, para attingir o seu objecto, emprega instrumentos materiaes, --- o escalpello, o microscopio, a balança; mas o pensamento não tem partes, em que se possa dividir, nem côr, que se veja, nem quantidade, que se pese. Quem viu as fracções da idea de \emph{justiça}, ou quem póde dizer se o pensamento é branco ou vermelho, se pesa dois ou tres grammas?\par Nem appellem os materialistas para as experiencias de \emph{J'owrens}. Este sabio constatou simplesmente que a integridade de certas partes do cerebro é indispensavel para o regular exercicio das faculdades intellectuaes. Mas, do facto de uma coisa ser indispensavel para a producção de um phenomeno, poderemos nós concluir que essa coisa é a \emph{causa} ou o orgão do phenomeno ? Não; a sciencia não acceita tal conclusão. Tambem os \emph{tuberculos quadrigemeos} são indispensaveis para o acto regular da visão, e todavia nenhum physiologista póde dizer que elles são o orgão central da visão. Do mesmo modo, se a \emph{physiologia} ensina que o cerebro é indispensavel para nós pensarmos, d'ahi não podemos concluir que o cerebro é o \emph{orgão} do pensamento.\par Coneluimos, pois, que nem a \emph{anatomia} nem a \emph{physiologia} podem demonstrar que o cerebro é o \emph{orgão} do pensamento. \emph{O} pensamento não é uma secreção, nem um movimento, nem outra qualquer funcção do cerebro, mas é uma coisa immaterial que escapa ao alcance da experiencia sensivel; e por isso, assim como só póde derivar de uma faculdade \emph{espiritual}, assim tambem só póde ser percebido por uma faculdade: \emph{espiritual}.\par Be o cerebro não é o \emph{orgão do pensamento}, é, todavia, uma sua \emph{condição}\emph{indispensavel}. \emph{Lelut}, seguindo as tradições da philosophia escholastica e christã, diz: O cerebro é o orgão das \emph{condições} do pensamento.\par Demonstraremos brevemente a verdade d'esta affirmação.\par O homem não é nem Anjo nem animal, mas participa de um e de outro, é um composto de materia e de espirito, é o \emph{animal racional}.\par Da essencia deriva a operação; esta é effeito d'aquella. Por isso entre uma e outra deve existir uma proporção, de modo que a operação reflicta os caracteres da essencia. Ora, a essencia do homem é composta de dois elementos, --- o \emph{material} e o \emph{espiritual}. Logo o pensamento deve manifestar esses dois elementos, --- o elemento \emph{material} e o elemento \emph{espiritual}.\par No pensamento devemos distinguir duas coisas: o \emph{principio}, de que elle deriva e que é a faculdade pensante, e o \emph{termo}, a que chega e que é o \emph{objecto} pensado. O pensamento não manifesta o \emph{elemento material} no seu \emph{principio}, porque o \emph{principio} do pensamento é a intelligencia, e esta é espiritual. Logo manifesta o elemento \emph{material} no seu \emph{termo}, emquanto o \emph{objecto} é apresentado por uma faculdade material e de um modo concreto.\par Com effeito, o homem, quando nasce, não possue idea alguma. A nossa intelligencia é uma \emph{toboa rasa}, em que tudo se póde escrever, mas nada está escripto. Só depois de alguns annos a intelligencia desenvolve-se, e os primeiros objectos, que conhece, são os corporeos, os materiaes. Mas esses objectos, para que sejam percebidos, devem unir-se á faculdade cognitiva. Como são materiaes, não podem unir-se por si mesmos, mas por uma sua imagem ou representação. Essa imagem, porém, é \emph{concreta} e \emph{singular}, porque representa a essencia do objecto, envolvida nas circumstancias ou condições individuaes de \emph{tempo}, \emph{logar}, \emph{figura}, \emph{cór}, \emph{grandeza}, etc.; e por isso só póde residir na \emph{imaginação}, que é uma \emph{faculdade sensitiva}, visto que na intelligencia só reside o \emph{abstracto}, o \emph{universal}.\par Todavia, a imagem \emph{concreta} e \emph{material}, se não existe na intelligencia, póde comtudo unir-se a esta faculdade. Porquanto, representando aquella imagem a \emph{essencia} do objecto envolvida nas condições materiaes, podemos, por meio da abstracção, prescindir das condições materiaes e só considerar á essencia. A essencia, abstrahindo das condições materiaes, é \emph{universal} e \emph{abstracta}, e, como tal, póde unir-se e une-se á intelligencia.\par Para o conhecimento intellectual, pois, é necessario o exercicio de duas faculdades --- da \emph{imaginação} e da \emph{intelligencia}. O exercicio da \emph{imaginação} é necessario para o da \emph{intelligencia}; porque a intelligencia só póde conhecer os objectos que lhe apresenta a imaginação. Um cego de nascimento comparará a côr vermelha ao som de uma flauta, não porque lhe falta a intelligencia, mas porque na sua imaginação não existe a imagem sensivel da côr vermelha, que elle nunca viu. --- O exercicio da imaginação não só precede, mas acompanha sempre o exercicio da intelligencia. A nossa intelligencia não pode pensar num objecto, sem que a imaginação forme uma imagem sensivel d'esse objecto (ao menos a imagem da \emph{palavra}), na qual de qualquer modo resplandeca o ideal, que contemplamos. Isto verifica-se não só quando o objecto da intelligencia á material mas tambem quando é immaterial, Não\par *podemos elevar-nos a \textsc{Deus}, senão por meio do \emph{visivel}; não podemos perceber o \emph{absoluto}, o \emph{infinito}, o \emph{necessario}, senão por meio do \emph{relativo}, do \emph{finito}, do \emph{contingente}.\par Ora a \emph{intelligencia} é uma faculdade \emph{inorganica} e reside na alma; a \emph{imaginação} é uma faculdade \emph{organica} e reside no cerebro. Logo, na operação intellectual, devem concorrer o cerebro com as suas vibrações e movimentos, e a intelligencia com as diversas formas de pensamento. Assim como, sem o exercicio da intelligencia, não existiria o \emph{principio} cognitivo; assim tambem, sem o exercicio do cerebro, faltaria o \emph{objecto}, que deve ser conhecido. O pensamento, pois, deriva \emph{exclusivamente} da intelligencia, mas exige, como \emph{condição indispensavel}, o concurso do cerebro.\par S. Thomaz tinha dito tudo n'estas poucas palavras: « Corpus requiritur ad actionem intellectus, \emph{non sicut organwm}, quo talis actio exerceatur, \emph{sed ratione objecti}. Phantasma enim comparatur ad intellectum sicut color ad visum. Sic autem indigere corpore non removet intellectum esse subsistenteni. Alioquin animal non esset aliquid subsistens, cum indigeat exterioribus sensibilibus ad sentiendum » (S. T. 1 p. q. 75, a. 2 ad 3).\par Dos principios expostos resultam legitimamente os seguintes corollarios, muito importantes não só para a psychologia mas tambem para a physiologia, --- corollarios, de que os materialistas erradamente se servem para combater a espiritualidade do pensamento e da alma humana.\par 1.º) Se o cerebro não estiver sufficientemente desenvolvido, pelo menos naquellas partes onde se acham localisados os sentidos internos, que concorrem para o acto intellectual e que são especialmente a \emph{memoria} e a \emph{imaginação}, não estará convenientemente desenvolvida a intelligencia.\par 2.º) Enfraquecidas as funcções do cerebro, como nas doenças e em muitos velhos, --- ou alteradas, como na loucura e na embriaguez, --- ou suspensas, como no somno profundo e no deliquio, ficará tambem enfraquecido, alterado ou suspenso o pensamento; porque, não funccionando, ou não funccionando regularmente a sensibilidade interna, a intelligencia ou não terá presente o objecto, ou não o terá presente como elle é na sua realidade. Por isso a loucura, e qualquer outra anomalia intellectual, se não póde dar-se na alma \emph{pianista} de \emph{Descartes}, nem na alma \emph{monada} de \emph{Leibnitz}, completamente isoladas do organismo e cheias de ideas innatas, bem pode dar-se na verdadeira alma humana, que, para adquirir os seus conhecimentos, precisa do concurso das faculdades sensitiyas.\par 3.º) Uma lesão, mais ou menos grave, naquellas partes do cerebro, onde se acham loealisados os sentidos internos necessarios para o exercicio do pensamento, produzirà tambem uma desordem, mais ou menos grave, na vida intellectual. Este corollario explica sufficientemente certos factos, que a physiologia verifica mas não sabe explicar, como é que a mais pequena lesão do cerebro produz, ás vezes, graves perturbações, ao passo que, outras vezes, graves lesões não produzem o menor desarranjo intellectual.\par 4.) Quanto mais intenso e prolongado fôr o pensamento, tanto maior será a actividade do cerebro, tanto mais elevado o grau da sua temperatura, tanto mais rapida a circulação do sangue, e tanto mais abundante a combustão da substancia cerebral. Com effeito, attesta-nos a consciencia que anossa alma, sem criar novas forças, póde fazer passar da potencia para o acto as que existem em nossos orgãos, e dirigil-as no seu desenvolvimento, provocando movimentos no organismo, por muitos modos e motivos. Ora, se a vontade actua nos musculos e nos orgãos locomotores, á intelligencia actua nos orgãos da memoria e da imaginação, para metter em movimento as imagens sensiveis, que são os materiaes das nossas abstracções, das ideas e de todos os actos intellectuaes. Mas a intelligencia não activa as operações sensitivas dos orgãos cerebraes, sem provocar um augmento nas operações nutritivas d'esses orgãos e por isso sem que o sangue afflua ao cerebro com as riquezas nutritivas, e accelere as reacções chimicas da nutrição.\par 5.º) Num trabalho intellectual muito intenso, sentimos cançaço no cerebro, e especialmente nos \emph{lóbos frontaes}. Dá-se tal phenomeno, não já porque as ideas estejam localisadas nalguma parte do cerebro (as ideas, não tendo figura nem dimensão, não podem residir num sujeito extenso), mas porque ha nesse trabalho phenomenos de ordem sensivel, que são as representações da \emph{imaginação}, localisadas na camada cerebral e provavelmente nos lóbos frontaes.\par \emph{M}\par Examinemos algumas affirmações e difficuldades dos materialistas.\par a) \emph{Moleschott} escreve : « Se analyse, \emph{pelos processos da chimica}, o cerebro humano, só encontro elementos materiaes e não o espirito; com a chimica decomponho \emph{tudo} em ammoniaco, agua e saes. Logo o principio do pensamento é materia e não o espirito. »\par Concedemos de bom grado que \emph{Moleschott}, \emph{pelos processos da chimica}, \emph{so} encontra elementos materiaes e não o espirito; pois o espirito não é composto de elementos materiaes. Se um naturalista podesse um dia encontrar na ponta do seu escalpello o espirito, já o espirito não seria espirito. Para se encontrar o espirito, devem empregar-se outros processos que não são os da chimica, devem empregar-se os processos da razão e então esta nos dirá que, realizando-se no homem phenomenos espirituaes, deve tambem existir nelle um principio proporcionado, que é o espirito. --- Concedemos tambem a \emph{Moleschott} que possa decompor tudo o que fôr \emph{chimicamente} composto; mas negamos que possa decompor o que não é nem chimica nem physicamente composto, como á a alma humana.\par b) \emph{Moleschott} continua: « Uma garrafa, contendo carbonato, ammoniaco, chloreto de potassio, phosphato de sodio, cal, etc., é \emph{dum modo ideal} o \emph{principio vital completo das plantas}, \emph{dos animaes e do proprio homem}. »\par Não é necessario que insistamos na refutação d'esses desvarios. A razão e a experiencia, como provámos, ensinam que o \emph{principio vital} não é materia, nem se identifica com as forças physicas e chimicas da materia. --- \emph{Cl}. \emph{Bernard} diz: « \emph{Poderia facilmente demonstrar que}, \emph{em physiologia}, o \emph{materialismo a nada guia e nada explica} » (\emph{La science experimentale}, pag. 961).\par c) \emph{Herzen} (\emph{Le cerveau et l'activité cérébrale}) \emph{e Büchner} (\emph{Force at matiere}) fazem as seguintes objecções: « A physiologia ensina --- que as operações intellectuaes, como a percepção, o juizo e o raciocinio, demandam tempo, --- que a modificação nas operações organicas importa uma modificação na consciencia e no pensamento, --- e que a razão enfraquece com o enfraquecimento dos orgãos, Ora estes factos provam que as operações intellectuaes derivam de uma faculdade organica n. :\par Repetimos que é a alma espiritual, e só ella, que pensa; mas que, para pensar, precisa do concurso das faculdades sensitivas ou organicas. Por isso no pensamento reflecte-se naturalmente toda e qualquer perturbação ou modificação da percepção sensitiva.\par Este principio explica os factos apontados.\par A percepção, o juizo e o raciocinio demandam tempo, não porque são operações materiaes e organicas, mas porque suppõem uma percepção sensitiva, que é impossivel sem uma operação organica e por isso sem uma mudança material \emph{successiva}, isto é, sem tempo.\par A modificação na consciencia e no pensamento; que deriva da modificação nas operações organicas, explica-se tambem pelo facto dos sentidos fornecerem á intelligencia a materia do pensamento.\par \emph{A} intelligencia enfraquece com o organismo, porque os orgãos enfraquecidos não fornecem a materia sufficiente para os actos intellectuaes.\par Todavia advertimos que a dependencia entre a faculdade intellectual e o cerebro não é tão absoluta, como se pensa. \emph{Longet} escreve: « Esta dependencia tão natural não é de tal modo absoluta, que se não dêm bastantes exemplos em sentido contrario. Encontram-se meninos debeis que espantam pela precocidade da razão e amplitude do seu espirito, velhos caducos e proximos do sepulchro que conservam intacto o seu juizo, a memoria, a força do genio e o ardor da coragem... Algumas vezes a loucura é acompanhada por uma lesão notavel dos centros nervosos; mas que deverá dizer-se dos casos, em que \emph{Esquirol}, \emph{Leuret} e outros auctores conscienciosos affirmam não ter encontrado vestigio de alteração no cerebro? Os annaes da sciencia fornecem um grande numero de factos perfeitamente observados de alterações profundas na massa cerebral, sem que durante a vida se haja manifestado a mais leve perturbação da inteligencia; têm-se visto porções de cerebro arrancadas, balas que atravessaram esse orgão de parte a parte sem desarranjo do espirito, ao passo que, ás vezes, leves transvasações de sangue, num ponto restricto, promovem a febre, excitam o delirio e produzem rapidamente a morte ». (\emph{Traité de phisiologie}, 1819, t. III, pag. 593). --- E basta uma só \emph{excepção} nessa dependencia entre a faculdade intellectual e o cerebro, para se demonstrar que o pensamento não deriva do cerebro; porque, se derivasse realmente, essa excepção não poderia ter logar.\par d) O dr. \emph{Luys}, na sua obra --- \emph{Le cerveau et ses fonctions}, diz que a psychologia não tem razão de ser, e que a physiologia e a anatomia explicam \emph{todo} o homem; pois que o estudo da composição do cerebro pertence á anatomia e o estudo das funcções d'esse orgão é objecto da physiologia.\par Concordamos com o dr. \emph{Luys} em que o estudo das funcções do cerebro é objecto da physiologia. Mas isto nada prova contra a nossa these, se o adversario não demonstrar que o pensamento é uma funeção do cerebro. E estamos convencidos de que nunca o demonstrará. E de que modo o poderia fazer? --- Pela razão? não: porque a razão diz que um phenomeno espiritual, como é o pensamento, não pbde ser funcção de um orgão material como éo cerebro. --- Pela experiencia? tambem não; porque nenhum physiologista viu funccionar o cerebro, pois que, para isso, seria mister arrancar as meninges, mas, se arrancarmos as meninges, o cerebro deixará de funccionar.\par e) Os materialistas objectum: Perturbado o cerebro, não ha memoria dos conhecimentos passados. Logo não é a alma espiritual, mas é a materia que recurda os pensamentos anteriores.\par Respondemos que, perturbado o cerebro, não ha memoria dos pensamentos passados, não porque seja a materia que pense e recorde, mas porque faltam as imagens sensiveis, e, faltando estas, falta á memoria o seu objecto. --- Pelo contrario, o phenomeno da memoria prova contra os adversarios. \emph{Claudio Bernard} diz: « O corpo humano é um composto de materias, que se renovam incessantemente. Todas as partes do corpo estão sujeitas a um perpetuo movimento de transformação. Cada dia perdeis alguma coisa do vosso ser physico, e pela alimentação substituis o que perdeis. De modo que, num espaço de oito annos, pouco mais ou menos, a yossa carne e os vossos ossos são substituidos por nova carne e por novos ossos, em consequencia das transformações successivas. A mão, com que escreveis hoje, não é inteiramente composta das mesmas moleculas que ha oito annos. A forma éa mesma, mas é uma nova substancia que a completa. O que digo da mão, dil.o-hei do cerebro. A abobada craneana não é occupada pela mesma materia que ha oito annos. Sendo assim, visto que tudo se muda no vosso cerebro, em oito annos, como é que vos lembraes de coisas que foram vistas ou ouvidas ha mais de oito annos? Se as coisas são, como pretendem certos physiologistas, alojadas, incorustadas nos lóbos do vosso cerebro, como é que sobrevivem a desapparição dos lóbos? Esses lóbos não são os mesmos que ha oito annos, e todavia a memoria guardou intacto o seu deposito. E' que no homem ha outra coisa que não é materia, e que é alguma coisa immaterial, permanente e independente da materia. Esta coisa é a ALMA ».}.

%%>>> t49

\pg{48} \tese{49}{Gonsequencias do materialismo.}
O \emph{materialismo} não só é falso nas doutrinas, mas é pernicioso nas consequencias; pois destroe a \emph{sciencia}, a \emph{arte}, \emph{á virtude}, a \emph{religião}, e á \emph{sociedade}.

a) Destroe a \emph{sciencia}. A sciencia é essencialmente dotada de \emph{unidade} e \emph{universalidade}. A \emph{unidade} é um effeito da \emph{generalisação}; \emph{a universalidade}, da \emph{abstracção}. Ora a intelligencia, se fosse \pg{50}  uma faculdade material ou organica, como ensina o materialiso, não poderia \emph{generalisar} nem \emph{abstrahir}; pois á operação não pode exceder a virtude do principio de que emana.

b) Destroe a \emph{arte}. \emph{A arte} não póde existir nem progredir, se o artista não formar um \emph{deal}, pelo qual transforme a materia. Ora, o artista não poderia formar o \emph{ideal}, se não possuisse \pg{52}  uma alma \emph{espiritual}, ou independente da materia; pois o \emph{ideal} forma-se na contemplação de ideas \emph{universaes}, que só podem ser concebidas por um principio espiritual.

\pg{49} \pg{51}
